\part{Unstable Chromatic Homotopy Theory}

\chapter{Chromatic localizations on spaces}

\section{More on chromatic localizations}

	We first give a short review on reflective localizations. 
	\begin{defn}
		For a family of morphisms $S$ in a category $\mathcal{C}$, we define an object $X$ to be $S$-local if for any $f:Y\to Z$ in $S$, we have $f^*:\Hom(Z,X)\to\Hom(Y,X)$ is an equivalence. We define $f:Y\to Z$ to be an $S$-equivalence if for all $X$ to be $S$-local, we have $f^*:\Hom(Y,X)\to\Hom(Z,X)$ an equivalence. Finally we define an object $X$ to be $S$-acyclic if $X\to *$ is an $S$-equivalence. 
	\end{defn}

	\begin{thm}\label{Refl_loc}
		There is a localization functor $L_S:\mathcal{C}\adj\mathcal{C}_S:\iota$ from $\mathcal{C}$ to the full-subcategory consisting $S$-local objects $\mathcal{C}_S$. Moreover, this localization functor exhibits an equivalence $\mathcal{C}_S\simeq\mathcal{C}[S^{-1}]$. By construction, we have a morphism $X\to L_S(X)$ which is an $S$-equivalence. 
	\end{thm}

	\begin{defn}\label{Null_loc_equiv}
		For any object $A$, we define the localization functor $P_A$ to be the reflective localization associated to the set consisting of a single map $A\to *$. 

		We say an object $X$ is $A$-null if we have $\Map(*,X)\to\Map(A,X)$ is an equivalence. This, equivalently, means $X$ is $S$-local where $S$ consists of the single morphism $A\to *$. 

		We say an object $Y$ is $A$-periodic if for all $A$-null spectra $X$, the inclusion $\Map(*,X)\to\Map(Y,X)$ is an equivalence. This means the morphism $Y$ is $S$-acyclic where $S$ is the same as in the above. 

		In particular, the above definitions apply to the category $\Sp$ or $\An$. 
	\end{defn}

	\begin{remark}
		The theory of Bousfield localizations is complete contained in reflective localizations. Note that for a spectra, $E$-equivalences ($E$-local objects) are simply the $S$-equivalences ($S$-local objects) in which $S$ is the set of morphisms $X\to *$ where $X$ runs through all $E$-acyclic spectra. 
	\end{remark}

	\begin{lemma}\label{Null_finite_smashing}
		For any spectrum $A$, $P_A$ commutes with finite colimits, hence for finite $F$, we have $P_A(X)\wedge F\simeq P_A(X\wedge F)$. 
	\end{lemma}
	\begin{proof}
		Similar to the proof of \cref{Loc_finite_smashing}. 
	\end{proof}

	\begin{lemma}\label{Null_prop_smash}
		For $E\wedge A\simeq*$, we have a natural equivalence $E\wedge X\to E\wedge P_A(X)$ for all $X$. 
	\end{lemma}
	\begin{proof}
		From the construction of $P_A$ (which I won't do here). (Yes, absolutely not). 
	\end{proof}

	\begin{defn}
		We define $F(n)$ to be a spectrum of type $n$. This is well-defined up to Bousfield equivalence. 
	\end{defn}

	\begin{lemma}\label{Loc_equiv_lem_1}
		For a spectrum $X$ and a self map $f:\Sigma^dX\to X$, we have $\langle X\rangle=\langle X[f^{-1}]\rangle\wedge\langle X/\Sigma^dX\rangle$. 
	\end{lemma}
	\begin{proof}
		For any spectrum $Y$, $Y\wedge X\simeq*$ obviously implies $Y\wedge X[f^{-1}]\simeq*$ and $Y\wedge X/\Sigma^dX\simeq*$. For the converse, if $Y\wedge X/\Sigma^dX=0$, then we have a natural equivalence $Y\wedge f:Y\wedge\Sigma^dX\to Y\wedge X$. Thus every map in the diagram $Y\wedge X\to Y\wedge\Sigma^{-d}X\to\cdots$ is an equivalence, hence induces an equivalence $Y\wedge X\to Y\wedge X[f^{-1}]$. Thus $Y\wedge X\simeq*$. 
	\end{proof}

	\begin{lemma}\label{Loc_equiv_lem_2}
		We have an orthogonal decomposition for any choice of $F(n)$
		\[ \langle S\rangle=\langle F(0)[v_0^{-1}]\rangle\vee\cdots\vee\langle F(n)[v_n^{-1}]\rangle\vee\langle F(n+1)\rangle \]
	\end{lemma}
	\begin{proof}
		By \cref{Loc_equiv_lem_1}, we have the decomposition. Then, to prove that it is orthogonal, it suffices to prove for $i<j$, $F(i)[v_i^{-1}]\wedge F(j)$ is contractible, so that $F(i)[v_i^{-1}]\wedge F(j)[v_j^{-1}]\simeq*$. However, supposing $v_i$ has degree $d$, this follows from that $v_i\wedge\id:\Sigma^dF(i)\wedge F(j)\to F(i)\wedge F(j)$ is zero on all $K(m)$, hence is nilpotent. 
	\end{proof}

	\begin{defn}
		We define $T(n)$ to be a spectrum of the form $X[f^{-1}]$, where $X$ is of type $n$ and $f$ is a $v_n$-self map of $X$. 
	\end{defn}

	\begin{prop}\label{Bousfield_class_telescope}
		Such $T(n)$ is well-defined up to Bousfield equivalence. 
	\end{prop}
	\begin{proof}
		For a type $n$ spectrum $X$ with a $v_n$-self map $f$, consider the full subcategory $\mathcal{D}_X$ of finite $p$-local spectra spanned by spectra $Y$ admitting a $v_n$-self map $g$ satisfying $X[f^{-1}]\wedge Z\simeq*\implies Y[g^{-1}]\wedge Z\simeq *$. 

		We then prove that $\mathcal{D}_X$ is a thick subcategory of finite $p$-local spectra. Now, $\mathcal{D}_X$ is obviously closed under retracts, since for $Y_1$ and $Y_2$ with morphisms $Y_1\to Y_2\to Y_1$, we may assume the $v_n$-self maps of $Y_i$ agrees with the morphisms. Then it follows that $Y_1[g_1^{-1}]\wedge Z$ is a retract of $Y_2[g_2^{-1}]\wedge Z$, hence is trivial. Then, for a fiber sequence $Y_1\to Y_2\to Y_3\to Y_4$ with $Y_2$ and $Y_3$ in $\mathcal{D}_X$, we may again assume the $v_n$-self maps of $Y_i$ agrees with this fiber sequence. Thus we have $Y_1[g_1^{-1}]\wedge Z$ and $Y_4[g_4^{-1}]\wedge Z$ are fibers or cofibers of amorphism between trivial spectra, hence trivial. 

		Finally, by \cref{Thick_subcat_thm} and that $\mathcal{D}_X$ contains $X$, we deduce that $\mathcal{D}_X\simeq\mathcal{C}_{\geq n}$. Therefore $T(n)$ is well-defined. 
	\end{proof}
	\begin{coro}\label{Tel_equiv_T_loc}
		We have $L_n^f\simeq L_{T(0)\vee T(1)\vee\cdots\vee T(n)}$. 
	\end{coro}
	\begin{proof}
		Recall that by \cref{Smash_loc_ker_type}, we have $L_n^f$ is the smashing localization having the minimal kernel with $\Sp_{(p)}^{\text{fin}}\cap\ker L_{E(n)}=\mathcal{C}_{\geq n+1}$. Then it suffices to prove $L_{T(0)\vee\cdots\vee T(n)}$ has the same property. Firstly, notice that by \cref{Loc_equiv_lem_2}, we have $\langle S\rangle=\langle T(0)\rangle\vee\cdots\langle T(n)\rangle\vee\langle F(n+1)\rangle$. Assuming for a spectrum $E$, we have $\Sp_{(p)}^{\text{fin}}\cap\ker L_{E(n)}\simeq\mathcal{C}_{\geq n+1}$, then smashing $E$ with the above equality, we have 
		\[ \langle E\rangle=\langle E\wedge T(0)\rangle\vee\cdots\langle E\wedge T(n)\rangle \]

		Then, it suffices to show that $L_{T(0)\vee\cdots\vee T(n)}$ is smashing. Notice that we have $\Sp_{(p)}^{\text{fin}}\cap\ker L_E=\Sp_{(p)}^{\text{fin}}\cap\ker L_{L_E(S)}$ for any spectrum $E$. Therefore we may apply $E=L_{T(0)\vee\cdots\vee T(n)}(S)$ to the above equality. In particular, we have 
		\[ \langle L_{T(0)\vee\cdots\vee T(n)}(S)\rangle\leq\langle T(0)\vee\cdots\vee T(n)\rangle \]

		However, since each $T(0)\vee\cdots\vee T(n)$-local spectrum is a module over $L_{T(0)\vee\cdots\vee T(n)}(S)$, $T(0)\vee\cdots\vee T(n)$-local implies $L_{T(0)\vee\cdots\vee T(n)}(S)$-local. Therefore, they have the same Bousfield class and hence $L_{T(0)\vee\cdots\vee T(n)}$ is smashing. 
	\end{proof}

	\begin{thm}\label{Tel_loc_equiv_F_null}
		We have $L_n^f\simeq P_{F(n+1)}$. 
	\end{thm}

	\begin{proof}
		By \cref{Loc_equiv_lem_2}, $(T(0)\vee\cdots\vee T(n))\wedge F(n+1)\simeq*$, thus $F(n+1)$ is $T(0)\vee\cdots\vee T(n)$-acyclic. Thus for all $X$, $\Map(F(n+1),L_n^f(X))\simeq*$. By the universal property of $P_{F(n+1)}$, this gives a natural transformation $\id\to P_{F(n+1)}\to L_n^f$. Thus to show the natural transformation is an equivalence, it suffices to show that for all $X$, $P_{F(n+1)}(X)$ is $T(0)\vee\cdots\vee T(n)$-local. To show this, let $Z$ be a $T(0)\vee\cdots\vee T(n)$-acyclic spectrum. We show that any map $Z\to P_{F(n+1)}(X)$ is null. It suffices to show that $P_{F(n+1)}(Z)$ is contractible. 

		By \cref{Loc_equiv_lem_2}, we have $T(i)\wedge F(n+1)\simeq*$ for $i\leq n$ and all possible choices of $T$ and $F$. Thus by \cref{Null_prop_smash}, we have $T(i)\wedge Z\to T(i)\wedge P_{F(n+1)}(Z)$ an equivalence. Since $Z$ is $T(0)\vee\cdots\vee T(n)$-acyclic, it is $T(i)$-acyclic. Therefore $T(i)\wedge P_{F(n+1)}(Z)\simeq*$. Then by definition of nullification, we have $\dual F(n+1)\wedge P_{F(n+1)}(Z)\simeq\Map(F(n+1),P_{F(n+1)}(Z))\simeq*$. But since $\langle \dual F(n+1)\rangle\simeq\langle F(n+1)\rangle$ and the equality in \cref{Loc_equiv_lem_2}, we actually see that $P_{F(n+1)}(Z)$ is $S$-null, hence $P_{F(n+1)}(Z)\simeq*$. 
	\end{proof}

\section{Localizations in the category $\An_*$}

	In this section, we apply the localizations introduced in the previous section to the category $\An_*$. The localizations in $\An_*$ is mostly the same as in $\Sp$. 
	A first remark here is that we would very much like to consider the localizations in the category $\An$ instead of $\An_*$ since it would be much more convenient to define the localizations (see the first remark below). However, there is no space $X$ with $\Sigma_+^\infty X$ a type-$n$ spectrum since it must contain a direct summand of the sphere spectrum $S$. On the otherhand, there is clearly a pointed space $Y$ such that $\Sigma^\infty Y$ is of type $n$ since we may consider the cellular structure of $F(n)$, and since it is finite, it must come from a finite pointed CW complex. 
	However, when we pass from the stable case to the unstable case, another problem we may encounter is that spectra representing the same Bousfield class (i.e. different spectra of type $n$), may end up having different Bousfield classes in $\An_*$ (if we regard $F(n)$ as a space with its suspension spectrum type $n$). This would be partially resolved by Bousfield's theorem (\cref{Bousfield_thm}), which mainly states that the unstable Bousfield classes of finite type $n$-spaces admitting certain properties is determined by the connectivity of the space. 

	\begin{const}
		In the category $\An_*$, the notion of $S$-local spaces and $S$-equivalences must be slightly modified. For pointed spaces $A$ and $B$, we let $\Hom(A,B)$ denote the mapping space of unpointed maps and $\Hom_*(A,B)$ denote the mapping space of pointed maps. Then we define the $S$-local spaces and $S$-equivalences in $\An_*$ to be respectively the objects $X$ such that $\Hom(B,X)\to\Hom(A,X)$ is an equivalence for all $A\to B\in S$ and the morphisms $Y\to Z$ such that $\Hom(Z,X)\to\Hom(Y,X)$ is an equivalence for all $S$-local space $X$. 
	\end{const}

	\begin{remark}
		For a self map $f:\Sigma^dA\to A$, the local objects associated to $f$ can be viewed as $Y$ such that $f_*:[A,Y]\to[\Sigma^dA,Y]$. This, in some sense, makes $L_f$ somewhat of a periodization. 
	\end{remark}

	We first present a few direct examples on localizations. 

	\begin{exam}\label{Loc_exam}
		Let $A=S^n$. Then a space is $A$-null iff it is $(n-1)$-truncated, and is $A$-periodic iff it is $n$-connective. 

		let $A=S^1/p$, then a simply connected $p$-local space is $A$-null iff it is rational, and is $A$-periodic iff its homotopy groups are torsion. 
	\end{exam}

	We might first notice that in the definition above, the choice of the $\Hom$ not preserving the base points is not that essential since we may always choose a base point. 
	
	\begin{lemma}\label{Basept_lemma}
		Let $S$ be a set of morphisms in $\An_*$ consisting of morphisms of connected spaces. Then a space $X$ is $S$-local iff for every choice of base point in $X$, the map $\Hom_*(B,X)\to\Hom_*(A,X)$ is an equivalence for all $A\to B$ in $S$. 
	\end{lemma}

	\begin{lemma}\label{Fibration_lemma}
		Let $S$ be a set of morphisms in $\An_*$. For a morphism of fibrations $f:(p_1:X_1\to B)\to(p_2:X_2\to B)$ inducing an $S$-equivalence (for some fixed set of morphisms $S$) for each $f_*:p_1^{-1}(b)\to p_2^{-1}(b)$, we have $f:X_1\to X_2$ is an $S$-equivalence. 
	\end{lemma}
	\begin{proof}
		Notice that we have natural identifications $X_i\simeq\ilim_{b\in B}p_i^{-1}(b)$, then the theorem follows by the fact that $S$-equivalences are preserved under colimits. 
	\end{proof}

	\begin{lemma}\label{Quotient_lemma}
		Suppose we have a H-group homomorphism $G_1\to G_2$ and a map of principal bundles. 
		% https://q.uiver.app/#q=WzAsNixbMCwwLCJHXzEiXSxbMSwwLCJYXzEiXSxbMiwwLCJCXzEiXSxbMCwxLCJHXzIiXSxbMSwxLCJYXzIiXSxbMiwxLCJCXzIiXSxbMCwxXSxbMCwzXSxbMSw0XSxbMSwyXSxbMiw1XSxbMyw0XSxbNCw1XV0=
		\[\begin{tikzcd}
			{G_1} & {X_1} & {B_1} \\
			{G_2} & {X_2} & {B_2}
			\arrow[from=1-1, to=1-2]
			\arrow[from=1-1, to=2-1]
			\arrow[from=1-2, to=1-3]
			\arrow[from=1-2, to=2-2]
			\arrow[from=1-3, to=2-3]
			\arrow[from=2-1, to=2-2]
			\arrow[from=2-2, to=2-3]
		\end{tikzcd}\]

		Moreover, the map $X_1\to X_2$ is required to be $G_1$-equivariant. If $G_1\to G_2$ and $X_1\to X_2$ are $S$-equivalences, then we have $B_1\to B_2$ an $S$-equivalence. 
	\end{lemma}
	\begin{proof}
		Simply write $B_i$ as the geometric realization of the simplicial object $X_1\times(G_1)^n$. 
	\end{proof}

	\begin{coro}\label{Quotient_lemma_coro_loopspace}
		If $X\to Y$ is a map of $n$-connective spaces having the property that $\Omega^nX\to\Omega^nY$ is an $S$-equivalence, then $X\to Y$ is an $S$-equivalence. 
	\end{coro}
	\begin{proof}
		Directly consider the path fibrations of $\Omega^{k-1}X$ and $\Omega^{k-1}Y$ and the corollary follows by induction. 
	\end{proof}

	\begin{lemma}\label{Smashing_lemma}
		If $Y$ is $S$-local, then $Y$ is $X\wedge S$-local for all spaces $X$, where $X\wedge S$ is the set of morphisms consisting of $X$ smashed with the morphisms in $S$. 
	\end{lemma}
	\begin{proof}
		Let $f:A\to B$ be a morphism in $S$. Then the space $X\wedge A$ is the homotopy cofiber of the map $X\vee A\inj X\times A$. This induces the space $\Map(X\wedge A,Y)$ as the homotopy pullback
		% https://q.uiver.app/#q=WzAsNCxbMCwwLCJcXEhvbShYXFx3ZWRnZSBBLFkpIl0sWzAsMSwiWSJdLFsxLDAsIlxcSG9tKFhcXHRpbWVzIEEsWSkiXSxbMSwxLCJcXEhvbShYXFx2ZWUgQSxZKSJdLFswLDJdLFsyLDNdLFswLDFdLFsxLDNdLFswLDMsIiIsMSx7InN0eWxlIjp7Im5hbWUiOiJjb3JuZXItaW52ZXJzZSJ9fV1d
		\[\begin{tikzcd}
			{\Hom(X\wedge A,Y)} & {\Hom(X\times A,Y)} \\
			Y & {\Hom(X\vee A,Y)}
			\arrow[from=1-1, to=1-2]
			\arrow[from=1-1, to=2-1]
			\arrow["\ulcorner"{anchor=center, pos=0.125}, draw=none, from=1-1, to=2-2]
			\arrow[from=1-2, to=2-2]
			\arrow[from=2-1, to=2-2]
		\end{tikzcd}\]

		Moreover $\Hom(X\vee A,Y)$ fits in a homotopy pullback square 
		% https://q.uiver.app/#q=WzAsNCxbMCwwLCJcXEhvbShYXFx2ZWUgQSxZKSJdLFswLDEsIlkiXSxbMSwwLCJcXEhvbShYLFkpXFx0aW1lc1xcSG9tKEEsWSkiXSxbMSwxLCJZXFx0aW1lcyBZIl0sWzAsMl0sWzIsM10sWzAsMV0sWzEsM11d
		\[\begin{tikzcd}
			{\Hom(X\vee A,Y)} & {\Hom(X,Y)\times\Hom(A,Y)} \\
			Y & {Y\times Y}
			\arrow[from=1-1, to=1-2]
			\arrow[from=1-1, to=2-1]
			\arrow[from=1-2, to=2-2]
			\arrow[from=2-1, to=2-2]
		\end{tikzcd}\]

		Thus it suffices to show that $\Hom(X\times B,Y)\to\Hom(X\times A,Y)$ is a weak equivalence. Rewriting this in terms of $\Hom(X,\Hom(B,Y))\to\Hom(X,\Hom(A,Y))$ gives the proof immediately. 
	\end{proof}

	The lemmas above gives rise to many corollaries on localizations, which we will list below. 

	\begin{thm}\label{Elementary_prop_loc}
		In the following list we assume $S$ and $T$ to be sets of morphisms in $\An_*$ and $X$ a pointed space. 
		\begin{enumerate}
			\item If $S$ consists of morphisms of connected spaces, then $X$ is $S$-local iff every path-component of $X$ is $S$-local (obvious). 
			\item Suppose $S$ consists of morphism of connected spaces, then $L_SX$ is connected implies that $X$ is connected (follows from the construction). 
			\item $L_S$ preserves finite products (follows from \cref{Smashing_lemma} and some theory on exponential ideals). 
			\item $L_S$ preserves colimits (follows from the fact that $L_S$ is a left adjoint). 
			\item A space is $S$-local and $T$-local iff it is $S\vee T$-local, where $S\vee T$ is the set of morphisms of the form $f\vee g$ where $f\in S$, $g\in T$ (direct). 
			\item A space $X$ is $FS$-local iff for every choice of basepoint $GX$ is $S$-local, where $F:\An_*\adj\An_*:G$ is a pair of adjoint functors (follows from \cref{Basept_lemma}). 
			\item If $X$ is $S$-acyclic, then $\Sigma X$ is $\Sigma S$-acyclic (follows from the properties of left adjoints).
			\item Let $f:X\to Y$ be a map of spaces. Assume $Y$ is $S$-local, then $X$ is $S$-local implies that for every point $y$ of $Y$ the fiber $f^{-1}(y)$ is $S$-local (by the long exact sequence of homotopy groups of fiber sequences). 
			\item Let $f:X\to Y$ be a map of spaces. Assume $Y$ is $A$-null, then $X$ is $A$-null iff for every point $y$ of $Y$ the fiber $f^{-1}(y)$ is $A$-null (similar to the above argument). 
			\item If $A$ is $n+1$-connective and $X$ is $n$-truncated, then $X$ is $A$-null (by the properties of truncatedness and connectivity). 
			\item If $S$ consists of $n$-connective maps, then $X\to L_SX$ is $n$-connective (by the construction of $L_SX$). 
		\end{enumerate}
	\end{thm}

	Then we introduce some less trivial facts on localizations. 

	\begin{lemma}\label{Per_implies_acyc}
		Let $S$ be a set of morphisms in $\An_*$, and $W$ is $S$-acyclic. Then for any space $X$, $X$ is $W$-periodic implies that $X$ is $S$-acyclic. In particular, if $W$ is $S$ acyclic having $\pi_0(W)\neq*$, then for all $X$ we have $X$ is $S$-acyclic. 
	\end{lemma}
	\begin{proof}
		We prove that if $Y$ is $S$-local, then $Y$ is $W$-null. Since $L_SW\simeq*$, by the universal property of the localization the inclusion $Y\inj\Hom(W,Y)$ is an equivalence, hence $Y$ is $W$-null. For the second assertion we notice $S^0$ is a retract of $W$, hence $S^0$ is $S$-acyclic. Thus any $S^0$-periodic space (hence every space) is $S$-acyclic. 
	\end{proof}
	
	\begin{lemma}\label{Susp_k_conn_acyc}
		Let $S$ be a set of morphisms in $\An_*$. If $L_SX\simeq*$, then for any $k$-connective space $Y$, we have $L_{\Sigma^kS}(X\wedge Y)\simeq*$. 
	\end{lemma}
	\begin{proof}
		Firstly the case where $Y\simeq S^n$ is true by entries (4) and (7) in \cref{Elementary_prop_loc}. Then we consider the general case by proving $L_{\Sigma^kS}(X\wedge \tau_{\leq n}Y)\simeq*$ by an induction on $n$. The base case $n=k$ is trivial since $\tau_{\leq k}Y$ is a wedge of $S^k$, and hence is obvious. Then the induction step follows from the fact that we have a cofiber sequence 
		\[ X\wedge\tau_{\leq n-1}Y\to X\wedge\tau_{\leq n}Y\to X\wedge\bigvee S^n \]

		Since the first space and last space is $\Sigma^kS$-acyclic, so is the middle space since $L_S$ preserves cofiber sequences. 

		Finally by taking a colimit $X\wedge Y\simeq\ilim X\wedge\tau_{\leq n}Y$, $X\wedge Y$ is $\Sigma^kS$-acyclic. 
	\end{proof}

	\begin{coro}\label{Susp_k_conn_acyc_sp}
		If $L_SX\simeq*$, then for any $k$-connective spectrum $Z$, we have 
		\[ L_{\Sigma^kS}\Omega^\infty(X\wedge Z)\simeq* \]
	\end{coro}
	\begin{proof}
		By \cref{Per_implies_acyc} it suffices to consider the case where $X$ is connected, otherwise every space is $f$-acyclic. Then since $L_SX\simeq *$, we have $L_{\Sigma^\infty S}\Sigma^\infty X\simeq*$ by a similar argument to entry (7) in \cref{Elementary_prop_loc}, hence $L_SQX\simeq *$, where $QX=\Omega^\infty\Sigma^\infty X$. Therefore to prove $L_{\Sigma^kf}\Omega^\infty(X\wedge Z)\simeq*$, it suffices to prove $L_{\Sigma^{\infty+k}f}X\wedge Z\simeq*$, which is by a similar argument to \cref{Susp_k_conn_acyc} 
	\end{proof}

	\begin{coro}\label{Homology_EM_periodic}
		If $\tilde{H}_d(A,\mathbb{F}_p)\neq0$, then for $d'\geq d$, we have $K(\mathbb{Z}/p\mathbb{Z},d')$ is $A$-periodic. 
	\end{coro}
	\begin{proof}
		Firstly by letting $X=A$, $Z=H\mathbb{F}_p$, $k=0$ in \cref{Susp_k_conn_acyc_sp}, since $P_AA\simeq*$, we have $P_A\Omega^\infty(H\mathbb{F}_p\otimes A)\simeq*$. Then since $\tilde{H}_d(A,\mathbb{F}_p)=G\neq0$, $\Omega^\infty(H\mathbb{F}_p\otimes A)$ contains $K(G,d)$ as a component in the product. Hence $K(G,d)$ is a retract of $\Omega^\infty(H\mathbb{F}_p\otimes A)$, and hence $P_AK(G,d)\simeq*$. Moreover, we have $G$ an $\mathbb{F}_p$-vector space, so $P_AK(\mathbb{F}_p,d)\simeq*$. Finally by \cref{Quotient_lemma_coro_loopspace} we have immediately $P_AK(\mathbb{F}_p,d')\simeq*$, since $\Omega^{d'-d}K(\mathbb{F}_p,d')\simeq K(\mathbb{F}_p,d)$ and $\Omega^{d'-d}*\simeq*$. 
	\end{proof}

	\begin{lemma}\label{Fib_nullification_periodic}
		Let $A$ be a space and $X$ be any space. Then consider the map $X\to P_AX$, and let $F$ be the fiber of the map at any point $p\in P_AX$. Then we have $P_AF\simeq*$. 
	\end{lemma}
	\begin{proof}
		By Lurie's straightening and unstraightening construction, the fibration $f:X\to P_AX$ is equivalently a functor $\St_f:P_AX\to\An$, where $f^{-1}(p)\simeq\St_f(p)$. Then we consider the composition $P_A\circ\St_f:X\to\An\to\An$, then this is equivalent to another map to $P_AX$, say $\overline{X}\to P_AX$. Note that the fiber of this map at $p$ is $P_AF$, and we have a morphism of fibrations 
		% https://q.uiver.app/#q=WzAsNixbMCwwLCJGIl0sWzEsMCwiWCJdLFsyLDAsIlBfQVgiXSxbMCwxLCJQX0FGIl0sWzEsMSwiXFxvdmVybGluZXtYfSJdLFsyLDEsIlBfQVgiXSxbMCwzXSxbMSw0XSxbMiw1XSxbMCwxXSxbMSwyXSxbMyw0XSxbNCw1XSxbMiw0LCIiLDEseyJzdHlsZSI6eyJib2R5Ijp7Im5hbWUiOiJkYXNoZWQifX19XV0=
		\[\begin{tikzcd}[ampersand replacement=\&]
			F \& X \& {P_AX} \\
			{P_AF} \& {\overline{X}} \& {P_AX}
			\arrow[from=1-1, to=1-2]
			\arrow[from=1-1, to=2-1]
			\arrow[from=1-2, to=1-3]
			\arrow[from=1-2, to=2-2]
			\arrow[dashed, from=1-3, to=2-2]
			\arrow[from=1-3, to=2-3]
			\arrow[from=2-1, to=2-2]
			\arrow[from=2-2, to=2-3]
		\end{tikzcd}\]

		And since the bottom fibration is of base and every fiber a $W$-null space, hence so is $\overline{X}$ by entry (9) of \cref{Elementary_prop_loc}. Therefore by the universal property of nullification, we get the lift in the diagram. Thus the map $F\to P_AF\to\overline{X}$ is null, which, again by the universal property of nullification, implies that the map $P_AF\to\overline{X}$ is null-homotopic. Moreover, the lift provides a section of the bottom fibration, hence the bottom fibration is trivial, and hence $P_AF$ is itself contractible. 
	\end{proof}

	\begin{lemma}\label{Loc_prod_EM_space}
		Let $S$ be a set of morphisms in $\An_*$. If $X$ is a product of Eilenberg-Maclane spaces $K(G,n)$ with all $G$ Abelian, then $L_S(X)$ is also a product of Eilenberg-Maclane spaces. 
	\end{lemma}
	\begin{proof}
		We first recall from the entry (3) of \cref{Elementary_prop_loc} that $L_S$ preserves finite products. Then we prove that a space is a product of Eilenberg-Maclane spaces iff it is of the form $\Omega^\infty M$ for some $H\mathbb{Z}$-module $M$. Only if is obvious by the delooping construction. If is by the following observation: for a $H\mathbb{Z}$-module $M$, the composition $M\to H\mathbb{Z}\otimes M\to M$ where the first map is the inclusion $\mathbb{S}\otimes M\to H\mathbb{Z}\otimes M$ induced by the unit of $H\mathbb{Z}$, hence $M$ is a retract of $H\mathbb{Z}\otimes M$. Then choose a Moore spectrum $M_n$ such that $\pi_n(M_n\otimes H\mathbb{Z})\simeq\pi_n(M)$ and a comparison map $M_n\to M$ inducing the isomorphism on $\pi_n$, and the map $H\mathbb{Z}\otimes(\bigoplus M_n)\to H\mathbb{Z}\otimes M\to M$ induces an isomorphism on homotopy groups, and hence a splitting $M\simeq\bigoplus H\mathbb{Z}\otimes M_n$. 

		Thus it suffices to see that $L_S$ lifts to the stable case, that is, there is some $\Omega^\infty\tilde{L}_SM\simeq L_S\Omega^\infty M$ for all connective $H\mathbb{Z}$-modules $M$. This follows from the fact that $L_S$ preserves products, hence preseserves infinite loopspaces. 
	\end{proof}

	\begin{lemma}\label{Commute_loc_loopspace}
		Let $S$ be a set of morphisms in $\An_*$. If $X$ is connected, then $L_S\Omega X\simeq\Omega L_{\Sigma S}X$. 
	\end{lemma}
	\begin{proof}
		Firstly we have a map $L_S\Omega X\to\Omega L_{\Sigma S}X$ given by the following process. For a connected space $X$, we have the localization map $X\to L_{\Sigma S}X$. Then apply the loopspace functor we have $\Omega X\to\Omega L_{\Sigma S}X$. Then by the entry (6) in \cref{Elementary_prop_loc}, $\Omega L_{\Sigma S}X$ is $S$-local, hence we have a factorization $\Omega X\to L_S\Omega X\to\Omega L_{\Sigma S}X$. 

		We then construct its inverse. Since $X$ is connected, so is $L_{\Sigma S}X$. Since $f$ preserves products, the map $\Omega X\to L_S\Omega X$ is a map of loopspaces, and we pass to classifying spaces we have a map $X\to BL_S\Omega X$. Then again by entry (6), $BL_S\Omega X$ is $L_{\Sigma S}$-local, and hence we have a factorization $X\to L_{\Sigma S}X\to BL_S\Omega X$. Applying the loopspace functor we have a map $\Omega L_{\Sigma S}\to L_S\Omega X$, which we would show to be the inverse. 

		We now verify that the maps are indeed inverses to each other. Since we have a map $\Omega X\to\Omega L_{\Sigma S}X\to L_S\Omega X$ whose composition is the localization of $\Omega X$, by the universal property of localizations $L_S\Omega X$ is also the localization of $\Omega L_{\Sigma S}X$, hence the composition $L_S\Omega X\to\Omega L_{\Sigma S}X\to L_S\Omega X$ is the identity. 

		For the other composite, we deloop the maps $\Omega X\to L_S\Omega X\to\Omega L_{\Sigma S}X$ and obtain that $X\to BL_S\Omega X\to L_{\Sigma S}X$ is the localization map. Hence we also have the $L_{\Sigma S}X\to BL_S\Omega X\to L_{\Sigma S}X$ is the identity, and passing to loopspaces we obtain the result. 
	\end{proof}

	\begin{lemma}\label{Susp_loc_prod_EM}
		Let $S$ be a set of morphisms in $\An$ consisting of morphisms between connected spaces. Let $X$ be a connected $H$-space, then if $L_SX\simeq*$ then we have $L_{\Sigma S}X$ a product of Abelian Eilenberg-Maclane spaces. 
	\end{lemma}
	\begin{proof}
		Firstly notice that $L_{\Sigma S}X$ is a connected $H$-space since $L_{\Sigma S}$ preserves finite products and connectivity (by the entry (3) and (11) in \cref{Elementary_prop_loc}). Moreover, since everything is connected here, we may use the space of pointed maps. 

		We first notice that we have 
		\[ \Hom_*(\Sigma L_{\Sigma S}X,L_{\Sigma S}X)\simeq\Hom_*(L_{\Sigma S}X,\Omega L_{\Sigma S}X)\simeq\Hom_*(L_{\Sigma S}X,L_SX)\simeq* \]

		Thus we have $L_{\Sigma S}X$ is $\Sigma L_{\Sigma S}X$-null. Then we prove $L_{\Sigma S}X$ is a retract of $\mathrm{SP}(L_{\Sigma S}X)$ (the infinite symmetric product), which finishes the theorem since by the Dold-Thom theorem, the symmetric product is a product of Abelian Eilenberg-Maclane spaces and a retract of Abelian Eilenberg-Maclane spaces is a product of Abelian Eilenberg-Maclane spaces. To prove this, we prove the inclusion $L_{\Sigma S}X\inj\mathrm{SP}(L_{\Sigma S}X)$ induces an isomorphism 
		\[ \Hom_*(\mathrm{SP}(L_{\Sigma S}X),L_{\Sigma S}X)\simeq\Hom_*(L_{\Sigma S}X,L_{\Sigma S}X) \]

		In fact, we claim the following more general lemma. For a connected space $A$ and a connected $\Sigma A$-null $H$-space $X$, we have an equivalence $\Hom_*(\mathrm{SP}(A),X)\simeq\Hom_*(A,X)$. We first notice that since $X$ is $\Sigma A$-null, the space $\Hom_*(A,X)$ is discrete. Moreover $\pi_1(X)$ acts trivially on $\Hom_*(A,X)$ since $X$ is a connected $H$-space. We claim that under these conditions 
		\[ \Hom_*(A\times A,X)\simeq\Hom_*(A,X)\times\Hom_*(A,X) \]

		This is obvious since we have $\Hom(A\times A,X)\simeq\Hom(A,\Hom(A,X))$, and we can extract $\Hom_*(A,X)$ from the fiber of $\Hom(A,X)\to X$ by evaluating on the basepoint of $A$. Therefore we have a natural $\mathbb{E}_\infty$-group structure on $\Hom_*(A,X)$. 
		Finally noticing that by Dold-Thom theorem, we have $\mathrm{SP}(A)\simeq\Omega^\infty H\mathbb{Z}\otimes A$ since they are both isomorphic to $\prod K(H_n(A),n)$. Thus it suffices to prove that 
		\begin{align*}
			&\Hom_*(\Omega^\infty(H\mathbb{Z}\otimes A),X) \\
			\simeq&\Hom(H\mathbb{Z},B^\infty\Hom_*(A,X)) \\
			\simeq&\Hom_{\mathbb{E}_\infty}(\mathbb{Z},\Hom_*(A,X)) \\
			\simeq&\Hom_*(A,X)
		\end{align*}
		The only non-trivial part is the first equivalence, which can be proven by considering the following inclusion of full subcategories of $\mathbb{E}_\infty$-groups into the category $\Fun(\Fin_*^{\op},\An)$ (the so-called $\Gamma$-spaces). The first equivalence is then a standard fact by abstract nonsense. 
	\end{proof}

\section{The unstable Bousfield class and Bousfield theorem}

	As before, we may introduce the notion of Bousfield classes to the category $\An_*$. 
	
	\begin{defn}\label{Unstable_Bousfield_class}
		We say two pointed spaces $A$ and $B$ are unstably Bousfield equivalent if the $A$-null spaces coincide with the $B$-null spaces. We denote $\langle A\rangle$ by the set of all spaces unstably Bousfield equivalent to $A$. Moreover, we define $\langle A\rangle\leq\langle B\rangle$ if $X$ is $B$-null implies that $X$ is $A$-null. 
	\end{defn}

	\begin{coro}
		We have $\langle A\wedge B\rangle\leq\langle A\rangle$. 
	\end{coro}

	\begin{lemma}\label{Equiv_Bousfield_class_cond}
		The following are equivalent. 
		\begin{enumerate}
			\item $\langle A\rangle\leq\langle B\rangle$.
			\item Every $A$-periodic space is $B$-periodic. 
			\item $A$ is $B$-periodic. 
			\item The unit $X\to P_BX$ factors as $X\to P_AX\to P_BX$. 
			\item $P_BP_AX\simeq P_BX$. 
		\end{enumerate}
	\end{lemma}
	\begin{proof}
		Direct. 
	\end{proof}

	We then prove the following fundamental theorem on the unstable Bousfield class related to telescopic localizations. 

	\begin{cond}\label{Cond_type_loc}
		We consider the following condition on a space $A$. We suppose $A$ is a finite $p$-local space of the form $\Sigma B$ with $B$ connected. Moreover we require $A$ to have integral homology bounded $p$-power torsion (which along with the condition $A=\Sigma B$ implies that $A$ is simply connected). 
	\end{cond}
	
	\begin{thm}[Bousfield]\label{Bousfield_thm}
		Let $A$ and $A'$ be finite $p$-local spaces satisfying \cref{Cond_type_loc} and with their suspension spectrum having type $>0$ (that is, the reduced integral homology is torsion), then we have $\langle A\rangle\leq\langle A'\rangle$ iff $\type A\geq\type A'$ (we abbreviate the type of $\Sigma^\infty A$ by $\type A$), and $\conn A\geq\conn A'$. 
	\end{thm}

	We will need the following result to prove \cref{Bousfield_thm}

	\begin{thm}\label{Fib_null_susp_EM_tors}
		Let $A$ be a $n$-connective pointed space such that the reduced homology groups $\tilde{H}_*(A,\mathbb{Z})$ are $p$-power torsion and $H_n(A;\mathbb{F}_p)\neq0$, then we have for any $k\geq1$, $\fib(P_{\Sigma^{k+1}A}X\to P_{\Sigma^kA}X)$ is equivalent to an Eilenberg-Maclane space $K(G,n+k)$, where $G$ is a $p$-torsion Abelian group. 
	\end{thm}
	\begin{proof}
		By replacing $A$ by $\Sigma^{k-1}A$, it suffices to show that the fiber $F$ of $P_{\Sigma^2A}X\to P_{\Sigma A}X$ is an Eilenberg-Maclane space of the form $K(G,n+1)$ with $G$ Abelian and $p$-torsion. 

		By \cref{Fib_nullification_periodic} we have $P_{\Sigma A}F\simeq*$ since the nullification $P_{\Sigma A}F\simeq P_{\Sigma A}P_{\Sigma^2A}F$. Moreover we have $P_{\Sigma^2A}F\simeq F$ since the two spaces are already $\Sigma^2A$-null. Then we claim $P_{\Sigma^2A}F$ is an $H$-space, so that we can apply \cref{Susp_loc_prod_EM}. The structure of an $H$-space on $P_{\Sigma^2A}F$ originates from the fact that $P_{\Sigma^2A}(F\vee F)\simeq P_{\Sigma^2A}(F\times F)$. This is because 
		\[ \fib(F\vee F\to F\times F)\simeq S^1\wedge\Omega F\wedge\Omega F\simeq\Omega F\wedge Z \]
		where $Z$ is simply connected, and $\Omega F\wedge Z$ is $\Sigma^2A$-null since $P_{\Sigma A}F\simeq*\implies P_A\Omega F\simeq *$, and by \cref{Susp_k_conn_acyc} we have $P_{\Sigma^2A}(\Omega F\wedge Z)\simeq*$. Thus the codiagonal $P_{\Sigma^2A}(F\vee F)\to P_{\Sigma^2A}F$ can be lifted to a map $P_{\Sigma^2A}(F\times F)\to P_{\Sigma^2A}F$, and hence $P_{\Sigma^2A}F$ is an $H$-space. 

		Now that $F$ is a product of Abelian Eilenberg-Maclane spaces, we show that there is in fact only one Abelian Eilenberg-Maclane spaces in the product and that the group is $p$-torsion. By the entry (11) of \cref{Elementary_prop_loc}, and that $A$ is $n$-connective, we have $\Sigma A\to *$ is $n+1$-connective, hence so is $F$. 

		Next we show $\pi_m(F)$ has no $p$-torsion for all $m\geq n+2$. Notice that $K(\pi_m(F),m)$ is a retract of $F$, hence it is $\Sigma^2A$-null and $\Sigma A$-periodic. Moreover, by \cref{Homology_EM_periodic}, since 
		\[ \tilde{H}_{n+1}(\Sigma A,\mathbb{F}_p)=\tilde{H}_{n+2}(\Sigma^2A,\mathbb{F}_p)=\tilde{H}_n(A,\mathbb{F}_p)=0 \]
		we have $K(\mathbb{F}_p,m)$ is $\Sigma A$-periodic for $m\geq n+1$ and $\Sigma^2A$-periodic for $m\geq n+2$. Hence the map $K(\pi_m(F),m)\to K(\pi_m(F)/\pi_m(F)[p],m)$ has fiber $K(V,m)$, where $V$ is a $\mathbb{F}_p$-vector space, hence the map is an equivalence after applying $P_{\Sigma^2A}$ for $m\geq n+2$. Then since $K(\pi_m(F),m)$ is $\Sigma^2A$-null, the map is actually an equivalence for $m\geq n+2$. Therefore $\pi_m(F)$ has no $p$-torsion for $m\geq n+2$. 

		Then we show $\pi_m(F)$ satisfies $\pi_m(F)\otimes\mathbb{Z}[p^{-1}]=0$ for all $m\geq n+2$. Consider the fiber sequence 
		\[ K(\pi_m(F)\otimes\mathbb{F}_p,m-1)\to K(\pi_m(F),m)\xrightarrow{\times p}K(\pi_m(F),m) \]

		Thus since $K(\pi_m(F)\otimes\mathbb{F}_p,m-1)$ is $\Sigma A$-periodic, the multiplication by $p$ map 
		\[ K(\pi_m(F),m)\to K(\pi_m(F),m) \] 
		is an equivalence after applying $P_{\Sigma A}$. Passing to the limit, we have 
		\[ K(\pi_m(F),m)\to K(\pi_m(F)\otimes\mathbb{Z}[p^{-1}],m) \] 
		is an equivalence after applying $P_{\Sigma A}$. Therefore $K(\pi_m(F)\otimes\mathbb{Z}[p^{-1}],m)$ is $\Sigma A$-periodic since so is $K(\pi_m(F),m)$. However, $K(\pi_m(F)\otimes\mathbb{Z}[p^{-1}],m)$ is also $\Sigma A$-null by the following argument. Since $A$ has reduced homology consisting of $p$-power torsion, it (and hence $\Sigma A$) has trivial reduced $\mathbb{Z}[p^{-1}]$-homology. By the universal coefficient theorem, it has also trivial $\mathbb{Z}[p^{-1}]$-cohomology. Therefore by the representability of cohomology, we have 
		\[ \Hom(\Sigma A,K(\pi_m(F)\otimes\mathbb{Z}[p^{-1}],m))\simeq K(\pi_m(F)\otimes\mathbb{Z}[p^{-1}],m) \]

		Thus $K(\pi_m(F)\otimes\mathbb{Z}[p^{-1}],m)$ is both $\Sigma A$-null and $\Sigma A$-periodic, hence $\pi_m(F)\otimes\mathbb{Z}[p^{-1}]$ is trivial. 

		Now we have proven that $\pi_m(F)=0$ for all $m\geq n+2$ since they have no $p$-torsion and trivial after inverting $p$. It remains to show that $\pi_{n+1}(F)$ is $p$-torsion. Since $K(\pi_{n+1}(F),n+1)\to K(\pi_{n+1}(F)/\pi_{n+1}(F)[p],n+1)$ is an equivalence after applying $P_{\Sigma A}$, $K(\pi_{n+1}(F)/\pi_{n+1}(F)[p],n+1)$ is $\Sigma A$-periodic. Similar to the above argument, notice that $\pi_{n+1}(F)/\pi_{n+1}(F)[p]$ is $p$-torsion free, hence the space is $\Sigma A$-null. Hence $\pi_{n+1}(F)/\pi_{n+1}(F)[p]=0$, and hence $\pi_{n+1}(F)$ is $p$-torsion. 
	\end{proof}

	\begin{proof}
		We now prove \cref{Bousfield_thm}. We first state an auxillary lemma. We claim that if $\Sigma^\infty A$ is contained in the thick subcategory generated by $\Sigma^\infty A'$, then there is some $d$ such that $\langle\Sigma^dA\rangle\leq\langle A'\rangle$. We consider the full subcategory $\mathcal{C}$ of finite $p$-local spectra consisting of spectra of the form $\Sigma^{\infty+n}B$, where $n$ is any integer (may be negative) and $B$ is a finite $p$-local space such that $\langle B\rangle\leq\langle A'\rangle$. Then obviously this category contains $\Sigma^\infty A'$, and it suffices to prove that it is thick, so that it must contain $\Sigma^\infty A$, and hence a suspension of $A$ satisfies the desired property. 

		For a cofiber sequence of spectra $P\to Q\to R$ with two entries in $\mathcal{C}$, we wish to prove the third is also in $\mathcal{C}$. Now by definition it is closed under arbitrary suspension and loopspaces, so WLOG we assume $P$ and $Q$ are in $\mathcal{C}$ and prove $R$ is in $\mathcal{C}$. Again by definition of $\mathcal{C}$, we suppose there is a cofiber sequence of spaces $L\to M\to N$ such that 
		\[ P\to Q\to R\simeq\Sigma^{\infty+k}L\to\Sigma^{\infty+k}M\to\Sigma^{\infty+k}N \]
		moreover we may assume $\langle A'\rangle\geq\langle L\rangle$ and $\langle A'\rangle\geq\langle M\rangle$, or equivalently $L$ and $M$ are $A'$-periodic. Then since $P_{A'}$ is a left adjoint, it preserves cofiber sequences, hence $P_{A'}L\to P_{A'}M\to P_{A'}N$ is isomorphic to a cofiber sequence $*\to *\to P_{A'}N$, which forces $P_{A'}N$ to be trivial. 

		Then for a retract diagram $P\to Q\to P$ with $Q\in\mathcal{C}$, again assume 
		\[ P\to Q\to P\simeq\Sigma^{\infty+k}L\to\Sigma^{\infty+k}M\to\Sigma^{\infty+k}L \]
		with $M$ $A'$-periodic. Then $L$ is $A'$-periodic since $A'$-periodicity is closed under retracts. Hence we have proven the auxillary lemma. 

		We now prove that if $\langle A\rangle\leq\langle A'\rangle$ then $\type A\geq\type A'$ and $\conn A\geq\conn A'$. For connectivity, notice that since $A'$ is $\conn A'$-connective, it is $S^{\conn A'}$-periodic. Hence we have $\langle A'\rangle\leq\langle S^{\conn A'}\rangle$, hence $A$ is also $\conn A'$-connective by transitivity of the partial order. Thus $\conn A\geq\conn A'$. For the type, assume $\widetilde{K(n)}_*(A')=0$, since $K(n)$ is a field, equivalently we have $\widetilde{K(n)}^*(A')=0$, hence 
		\[ \pi_m\Hom(A',\Omega^\infty K(n))\simeq\pi_0\Hom(A',\Omega^{\infty+m}K(n))\simeq\widetilde{K(n)}^{-m}(A')\simeq* \]
		and thus $\Omega^{\infty+m}K(n)$ are $A'$-null for all $m$, and hence are $A$-null for all $m$ by assumption. Thus $\widetilde{K(n)}_*(A)=0$, and hence $\type A\geq\type A'$. 

		Finally it suffices to prove the converse. Since $\type A\geq\type A'$, by the thick subcategory theorem we clearly have $\Sigma^\infty A$ in the thick subcategory generated by $\Sigma^\infty A'$, hence there is some $d$ such that $\langle\Sigma^dA\rangle\leq\langle A'\rangle$. We will prove that $\langle\Sigma^{d-1}A\rangle\leq\langle A'\rangle$, so that by induction we will have $\langle A\rangle\leq\langle A'\rangle$. Hence we may assume that $d=1$. By the assumption, $P_{A'}(A)\simeq P_{A'}P_{\Sigma A}(A)$. Hence to show $\langle A\rangle\leq\langle A'\rangle$, it suffices to prove $P_{A'}P_{\Sigma A}(A)\simeq*$. But note that obviuously we have $P_A(A)\simeq*$, hence $P_{\Sigma A}(A)$ is the fiber of the map $P_{\Sigma A}(A)\to P_A(A)$, which, by \cref{Fib_null_susp_EM_tors}, has the form $K(G,d')$ where $G$ is $p$-power torsion. But this is clearly $A'$-periodic by \cref{Homology_EM_periodic}. Hence we have proven the theorem. 
	\end{proof}

\chapter{The Bousfield-Kuhn functor}

\section{Unstable telescopic localizations}

	With the preparations in the previous sections, we can finally begin the study of telescopic localizations on spaces. 

	Recall that the telescopic localization $L_n^f$ is given by killing all the $v_m$-maps for $m<n$ and inverting the $v_n$-map. In \cref{Tel_loc_equiv_F_null}, we find that $L_n^f\simeq P_{F(n+1)}$, which gives two possible methods to define telescopic localization on spaces. We first focus on one. We remark here that from now on, when we mention a type $n$ space, it will be implicitly be $p$-local for some fixed $p$. 

	\begin{defn}
		We fix integers $n$ and $d$. Suppose $d$ is sufficiently large so that there exists a space $F(n+1,d)$ of type $n+1$ and connectivity $d$ satisfying \cref{Cond_type_loc}. Then we define the category $L_n^f(\An_*)_{(p)}^{>d}$ to be the full subcategory of $\An_*$ consisting of the spaces that is $(d+1)$-connective, $p$-local and $F(n+1,d)$-null. By definition, this category admits all limits and colimits. 
	\end{defn}

	\begin{coro}
		We have a localization functor $P_{F(n+1,d)}:(\An_*)_{(p)}^{>d}\to L_n^f(\An_*)_{(p)}^{>d}$. 
	\end{coro}
	\begin{proof}
		The proof is mainly due to the definition of local objects. We only need to deal with the connectivity conditions here. Notice that $F(n+1,d)\to *$ is $d$-connective, so by the entry (11) of \cref{Elementary_prop_loc}, so is $X\to P_A(X)$. Now $X$ and $P_A(X)$ have isomorphic homotopy groups up to degree $d-1$, and we show $\pi_d(P_A(X))=0$. We have an exact sequence of homotopy groups 
		\[ \pi_d(X)\to\pi_d(P_A(X))\to\pi_{d-1}(X,P_A(X)) \] 
		with the first and last term zero, which forces $\pi_d(P_A(X))=0$. 
	\end{proof}

	The difference between the categories $L_n^f(\An_*)_{(p)}^{>d}$ when the type $n$ varies is analogous to the stable case, but we do encounter something new when $d$ varies. A good news is that this difference is simple enough. 

	\begin{prop}\label{Tel_loc_conn_var_rational}
		There is a fully faithful embedding $L_n^f(\An_*)_{(p)}^{>d+1}\inj L_n^f(\An_*)_{(p)}^{>d}$ which identifies the image as the spaces $X\in L_n^f(\An_*)_{(p)}^{>d}$ such that $(\pi_{d+1}X)_\mathbb{Q}$-vanishes. 
	\end{prop}
	\begin{proof}
		We consider the functor as the composition 
		\[ L_n^f(\An_*)_{(p)}^{>d+1}\inj(\An_*)_{(p)}^{>d}\xrightarrow{P_{F(n+1,d)}}L_n^f(\An_*)_{(p)}^{>d} \]
		which we will abuse notation to denote by $P_{F(n+1,d)}$. Then for $X\in L_n^f(\An_*)_{(p)}^{>d+1}$, the map $X\to P_{F(n+1,d)}X$ has fiber a $F(n+1,d)$-periodic space by the \cref{Fib_nullification_periodic}. Now, by \cref{Loc_exam}, recall that $S^1/p$-periodic $p$-local spaces have torsion homotopy groups. Then since $\Sigma S^1/p$ have type 1 and connectivity 2, we have $\langle\Sigma S^1/p\rangle\geq\langle F(n+1,d)\rangle$. Thus the fiber is also $\Sigma S^1/p$-periodic, hence trivial on rational homotopy groups. Thus $X\to P_{F(n+1,d)}X$ is an isomorphism on rational homotopy groups, hence we have the image of $P_{F(n+1,d)}$ contained in spaces where the $(d+1)$-th rational homotopy groups zero. 

		Then we consider the functor $\tau_{\geq d+2}:L_n^f(\An_*)_{(p)}^{>d}\to(\An_*)_{(p)}^{>d+1}$. For $X\in L_n^f(\An_*)_{(p)}^{>d}$, we have a fiber sequence $\tau_{\geq d+2}X\to X\to K(\pi_{d+1}(X),d+1)$. Since $X$ is $F(n+1,d)$-null, it is $\Sigma F(n+1,d)$-null by \cref{Smashing_lemma}. For $X$ with $(\pi_{n+1}(X))_\mathbb{Q}$ non-trivial, $K(\pi_{n+1}(X),d+1)$ is clearly $\Sigma F(n+1,d)$-null by the definition of spaces of type $n+1$. So when $(\pi_{n+1}(X))_\mathbb{Q}$ is non-trivial, $\tau_{\geq d+2}$ upgrades to a functor $L_n^f(\An_*)_{(p)}^{>d}\to L_n^f(\An_*)_{(p)}^{>d+1}$. 

		Finally we show that for $X$ with $(\pi_{n+1}(X))_\mathbb{Q}$-trivial, $\tau_{\geq d+2}$ is an inverse to the functor $P_{F(n+1,d)}$. Firstly for a space $Y\in L_n^f(\An_*)_{(p)}^{>d+1}$, we prove $Y\to P_{F(n+1,d)}(Y)$ induces an isomorphism on homotopy groups in degrees $>d+1$. Let $F$ be its fiber. Since $Y$ is $P_{\Sigma F(n+1,d)}$-local, $F$ is equivalently the fiber of $P_{\Sigma F(n+1,d)}(Y)\to P_{F(n+1,d)}(Y)$, hence has the form $K(G,d)$ by \cref{Fib_null_susp_EM_tors}. By delooping we have a fiber sequence $Y\to P_{F(n+1,d)}(Y)\to K(G,d+1)$, so $Y$ must be the $(d+2)$-connective cover of $P_{F(n+1,d)}(Y)$. 

		Then for a space $X\in L_n^f(\An_*)_{(p)}^{>d}$ with $(\pi_{d+1}(X))_\mathbb{Q}=0$, we prove the map $\tau_{\geq n+2}X\to X$ is a $F(n+1,d)$-nullification of $\tau_{\geq n+2}X$. We show that the fiber of this map is $F(n+1,d)$-periodic. Note that the fiber is the Eilenberg-Maclane space $K(\pi_{d+1}(X),d)$, but $\pi_{d+1}(X)$ is $p$-power torsion, and hence is periodic by \cref{Homology_EM_periodic}. 
	\end{proof}

	\begin{lemma}\label{Tel_loc_acyc_cond}
		Let $X$ be a $p$-local space. Then $X$ is $F(n+1,d)$-periodic iff the three following conditions are satisfied. 
		\begin{enumerate}
			\item $X$ is $d$-connective. 
			\item The homotopy group $\pi_d(X)$ is $p$-power torsion. 
			\item The connective cover $\tau_{\geq d+1}X$ is $\Sigma F(n+1,d)$-periodic. 
		\end{enumerate}
	\end{lemma}
	\begin{proof}
		If the three conditions are satisfied, then we have a fiber sequence $\tau_{\geq d+1}X\to X\to K(\pi_d(X),d)$, where $K(\pi_d(X),d)$ and $\tau_{\geq d+1}X$ are $F(n+1,d)$-periodic by \cref{Homology_EM_periodic} and definition, hence $X$ is $F(n+1,d)$-periodic. 

		Conversely, suppose $X$ is $F(n+1,d)$-periodic. By the entry (11) of \cref{Elementary_prop_loc} $X\to P_{F(n+1,d)}X\simeq*$ is $d$-connective. We have $P_{\Sigma F(n+1,d)}X$ equivalent to the fiber of $P_{\Sigma F(n+1,d)X}\to P_{F(n+1,d)}X\simeq*$, hence $P_{\Sigma F(n+1,d)}X\simeq K(G,d)$ for some $p$-power torsion Abelian group. Then consider the fiber sequence $F\to X\to P_{\Sigma F(n+1,d)}X\simeq K(G,d)$. Now $F$ is $\Sigma F(n+1,d)$-periodic, hence $(d+1)$-connective. This implies that $F$ is the $(d+1)$-connective cover of $X$, hence the second and third condition follows immediately. 
	\end{proof}

	\begin{prop}\label{Tel_loc_fin_lim}
		The functor $\An_*\to L_n^f(\An_*)_{(p)}^{>d}$ given by $X\mapsto P_{F(n+1,d)}(\tau_{\geq d+1}X_{(p)})$ is left exact (that is, preserves finite limits). We denote this functor by $L_{n,d}^f$. 
	\end{prop}
	\begin{proof}
		Obviously taking connective covers commutes with limits. Localization at $p$ is left exact since it firstly preserves the final object, and preserves pullbacks due to the Mayer-Vietoris sequence of homotopy groups of pullbacks, and the fact that $\mathbb{Z}_{(p)}$ is flat. Hence it now suffices to prove the functor $P_{F(n+1,d)}:(\An_*)_{(p)}^{>d}\to(\An_*)_{(p)}^{>d}$ is left exact. Again, it obviously preserves final objects, so it suffices to check that it preserves pullbacks. 

		Consider maps of $d$-connected pointed spaces $X_1\to X_0\from X_2$. We show the canonical map 
		\[ \tau_{\geq d+1}(X_1\times_{X_0}X_2)\to\tau_{\geq d+1}(P_{F(n+1,d)}X_1\times_{P_{F(n+1,d)}X_0}P_{F(n+1,d)}X_2) \]
		has $F(n+1,d)$-periodic fiber. We let $F$ denote this fiber. Obviously $F$ is $d$-connective, and since rationalization is left exact, $\pi_*(F)_\mathbb{Q}$ is trivial, hence the homotopy groups of $F$ is $p$-power torsion. Thus by \cref{Tel_loc_acyc_cond}, by induction, it suffices to find some $P_A$-local space $F'$ such that their $d'$-connective cover agree for $d'\gg d$. To do this, we let $Y_i$ be the fiber of the map $X_i\to P_{F(n+1,d)}X_i$, and consider $F'=\tau_{\geq d+1}Y_1\times_{\tau_{\geq d+1}Y_0}\tau_{\geq d+1}Y_2$. Now they have a common connctive cover $\tau_{\geq d+1}(Y_1\times_{Y_0}Y_2)$ by considering the long exact sequence of homotopy groups of a fiber sequence or a pullback square, hence by \cref{Tel_loc_acyc_cond} it suffices to prove $F'$ is $F(n+1,d)$-periodic. 

		Notice that 
		\[ P_{F(n+1,d)}\Omega\tau_{d+1}Y\simeq\Omega P_{\Sigma F(n+1,d)}\tau_{d+1}Y\simeq * \]
		where the first equivalence is \cref{Commute_loc_loopspace} the second equivalence is by \cref{Tel_loc_acyc_cond}. Thus $\Omega\tau_{d+1}Y_0$ is $F(n+1,d)$-periodic. Then we consider the fiber sequence $K\to\tau_{d+1}Y_1\to\tau_{d+1}Y_0$. This gives a fiber sequence $\Omega\tau_{d+1}Y_0\to K\to\tau_{d+1}Y_1$. Hence $K$ is $F(n+1,d)$-periodic since both $\Omega\tau_{d+1}Y_0$ and $\tau_{d+1}Y_1$ is. Now consider the fiber sequence $K\to\tau_{d+1}Y_1\times_{\tau_{d+1}Y_0}\tau_{d+1}Y_2\to\tau_{d+1}Y_2$, which proves that $F'$ is $F(n+1,d)$-periodic. This finishes the whole proof. 
	\end{proof}

	We now wish to investigate the properties preserved by $L_{n,d}^f$, which also leads to the other route to defining telescopic localizations in spaces. As before, we start from the stable setting. 

	\begin{defn}\label{Periodic_homotopy_grp_sp}
		For a spectrum $X$ and a finite $p$-local spectrum $F(n)$ of type $n$, we define the $v_n$-periodic homotopy groups of $X$ by 
		\[ v_n^{-1}\pi_d(X;F(n))\simeq\ilim(\pi_d\Map(F(n),X)\xrightarrow{v_n}\pi_{d+t}\Map(F(n),X)\to\cdots) \]
		where the colimit is given by precomposing with $v_n$. This construction only depends on the choice of $F(n)$ by \cref{Periodicity_thm_lem_5}. 
	\end{defn}

	\begin{prop}\label{T_n_loc_crit}
		A map of spectra $f:X\to Y$ is an $T(n)$-equivalence iff it induces an isomorphim on $v_n$-periodic homotopy groups for any choice of $F(n)$. 
	\end{prop}
	\begin{proof}
		We prove that we have the following identification of $v_n^{-1}\pi_*$ as a homology theory
		\[ v_n^{-1}\pi_*(X)\simeq DT(n)_*(X) \]
		where $DT(n)$ is the dual of the telescope of $F(n)$. Let $v_n$ be a $v_n$-self map of $F(n)$, then $Dv_n$ is a $v_n$-self map of $\dual F(n)$. Hence we may upgrade the construction of $v_n^{-1}\pi_d(X;F(n))$ to the colimit of spectrum 
		\[ \ilim(\Map(F(n),X)\to\Map(\Sigma^tF(n),X)\to\cdots) \]

		By definition of $p$-local finite spectra and duals, this is equivalently the colimit 
		\[ \ilim(\dual F(n)\otimes X\from\Sigma^{-t}\dual F(n)\otimes X\to\cdots) \]
		and by the definition of $T(m)$ in \cref{Telescope_const}, this is obviously $DT(m)\otimes X$. Hence we have the identification above. 
	\end{proof}

	We then turn to the unstable setting. 

	\begin{defn}\label{Periodic_homotopy_grp}
		For a space $X$, and a finite space of type $m$ equipped with a $v_n$-self map $v_n:\Sigma^tF(n)\to F(n)$, we define its $v_n$-periodic homotopy groups $v_n^{-1}\pi_d(X;F(n))$ by 
		\[ v_n^{-1}\pi_d(X;F(n))\simeq\ilim(\pi_d\Hom_*(F(n),X)\xrightarrow{v_n}\pi_{d+t}\Hom_*(F(n),X)\to\cdots) \]
		where the colimit is given by precomposing with $v_n$. This construction depends on the choice of $F(n)$, but not the choice of the $v_n$-self map (we will see this later). 
	\end{defn}

	\begin{defn}\label{Periodic_equiv}
		We define a map $f:X\to Y$ to be $v_n$-periodic equivalence if for all choices of base points and some $F(n)$, we have $v_n^{-1}\pi_*(X;F(n))\to v_n^{-1}\pi_*(Y;F(n))$ is an equivalence. 
	\end{defn}

	\begin{const}\label{Pre_Bousfield_Kuhn_fun}
		We may realize $v_n^{-1}\pi_d(X;F(n))$ as the homotopy groups of a spectra, since by definition it is periodic of period $t$. Explicitly, the spectra is constructed as follows. Let 
		\[ Y=\Hom_*(F(n),X)\to\Hom_*(\Sigma^tF(n),X)\to\Hom_*(\Sigma^{2t}F(n),X)\to\cdots \]

		Then clearly $\Omega^tY=Y$. Hence we define the spectrum $\Phi_{v_n}(X)$ by letting its $k$-th space be $\Omega^{Nt-k}Y$, where $N\gg0$ is any integer such that $Nt>k$. 
	\end{const}

	Not so obviously, the above construction does not depend on the choice of the $v_n$-self map by the following arguments. 

	\begin{lemma}\label{Pre_Bousfield_Kuhn_susp_power}
		Let $v_n:\Sigma^tF(n)\to F(n)$ be a $v_n$-self map of a finite space of type $n$. Then $\Phi_{\id_W\wedge v_n}(X)\simeq\Map(W,\Phi_{v_n}(X))$, $\Phi_{v_n^k}(X)\simeq\Phi_{v_n}(X)$. Here $\id_W\wedge v_n$ is the map $\id_W\wedge v_n:W\wedge\Sigma^tF(n)\to W\wedge F(n)$, where $W$ is a pointed space, and $v_n^k:\Sigma^{kt}F(n)\to F(n)$ be the map obtained by $v_n$ composed with itself by $k$ times. We remark here that by the Kunneth formula for the $K(n)$-homology, $W\wedge F(n)$ has always type $\geq n$, and when it has type $>n$, the $v_n$-self map $\id_W\wedge v_n$ is nilpotent, and $\Phi_{\id_W\wedge v_n}(X)\simeq0$. 
	\end{lemma}
	\begin{proof}
		Direct. 
	\end{proof}

	\begin{coro}\label{Pre_Bousfield_Kuhn_unique}
		For two different $v_n$-self maps $v_n:\Sigma^tF(n)\to F(n)$ and $v'_n:\Sigma^{t'}F(n)\to F(n)$, we have $\Phi_{v_n}\simeq\Phi_{v'_n}$. 
	\end{coro}
	\begin{proof}
		Recall from \cref{Periodicity_thm_lem_5} that we may assume $v_n$ and $v'_n$ are stably equivalent after raising to a suitable power. Thus $\Sigma^kv_n^a\simeq\Sigma^k(v'_n)^b$. By \cref{Pre_Bousfield_Kuhn_susp_power} we have $\Phi_{v_n}\simeq\Sigma^k\Phi_{\Sigma^kv_n^a}\simeq\Sigma^k\Phi_{\Sigma^k(v'_n)^b}\simeq\Phi_{v'_n}$. 
	\end{proof}

	Therefore by the uniqueness of $\Phi_{v_n}$, we will denote it by $\Phi_{F(n)}$ to show its independence of the choice of the $v_n$-self map. But do notice that the choice of the type $n$ space $F(n)$ still matters in the construction. 
	We also remark here that the construction $\Phi_{F(n)}(X)$ is also somewhat functorial in the choice of $F(n)$. Namely we have the following property. 
	
	\begin{lemma}\label{Pre_Bousfield_Kuhn_cofib}
		For a map of type $n$ spaces $F(n)\to F'(n)$, we have a functor $\Phi_{F'(n)}(X)\to\Phi_{F(n)}(X)$. Moreover, if we regard $F''(n)=\cof(F(n)\to F'(n))$ as a space equipped with a $v_n$-self map, then we have 
		\[ \Phi_{F''(n)}=\fib(\Phi_{F'(n)}(X)\to\Phi_{F(n)}(X)) \]
	\end{lemma}
	\begin{proof}
		Firstly the functorality is obtained by \cref{Periodicity_thm_lem_5}. The second property is by the fact that a cofiber of a map of $v_n$-self maps is a $v_n$-self map, and taking fibers commute with filtered colimits. 
	\end{proof}

	\begin{thm}\label{Periodic_equiv_indep_type_n}
		Let $f:X\to Y$ be a pointed map of pointed spaces. Then for any type $n$ spaces $F(n)$ and $F'(n)$ equipped with $v_n$-self maps, then $f$ induces an equivalence $\Phi_{F(n)}(X)\simeq\Phi_{F(n)}(Y)$ iff $\Phi_{F'(n)}(X)\simeq\Phi_{F'(n)}(Y)$. 
	\end{thm}
	\begin{proof}
		It suffices to prove that $X\to Y$ induces an isomorphism on $v_n$-periodic homotopy groups with coefficients in $F'(n)$. Recall that by the thick subcategory theorem, $\Sigma^\infty F'(n)$ is in the thick subcategory generated by $\Sigma^\infty F(n)$. Therefore $\Sigma^{\infty+N}F'(n)$ is generated by $\Sigma^\infty F(n)$ under retracts, suspensions and cofibers under finitely many steps for $N\gg0$. Therefore to prove the above, it suffices to prove the three special cases. 
		\begin{enumerate}
			\item $F'(n)$ is a retract of $F(n)$. 
			\item $F'(n)$ is a suspension of $F(n)$. 
			\item For another $F''(n)$ such that $\Phi_{F''(n)}(X)\simeq\Phi_{F''(n)}(Y)$ and a map $F''(n)\to F(n)$, $F'(n)=\cof(F''(n)\to F(n))$. 
		\end{enumerate}
		The first case is obvious since if this is the case then $v_n^{-1}\pi_*(X;F'(n))$ is a retract of \\
		$v_n^{-1}\pi_*(X;F(n))$, and a retract of an isomorphism is an isomorphism. The second case is by \cref{Pre_Bousfield_Kuhn_susp_power}. Finally the third case is by the fact that 
		\[ \Hom_*(F'(n),X)\simeq\fib(\Hom_*(F(n),X)\to\Hom_*(F''(n),X)) \]
		and that taking fibers commutes with filtered colimits. 
	\end{proof}

	Therefore we may define the $v_n$-periodic homotopy equivalences to the maps send by $\Phi_{F(n)}$ to equivalences for any type $n$ space $F(n)$. 

	\begin{thm}\label{Pre_Bousfield_Kuhn_T_n_loc}
		For any type $n$ space $F(n)$ equipped with a $v_n$-self map $v_n$, we have $\Phi_{F(n)}(X)$ is $T(n)$-local. 
	\end{thm}
	\begin{proof}
		Recall that we have from \cref{Tel_loc_equiv_F_null} that $L_n^f\simeq P_{F(n+1)}$ and from \cref{Tel_equiv_T_loc} that a space is $T(n)$-local iff it is $L_n^f$-local but not $L_{n-1}^f$-local. Therefore it suffices to see that $\Phi_{F(n)}(X)$ is $\Sigma^\infty F'(n+1)$-acyclic but not $\Sigma^\infty F'(n)$-acyclic for some type $n$ space $F'(n)$ and some type $n+1$ space $F'(n+1)$. The first fact follows from \cref{Pre_Bousfield_Kuhn_susp_power} since $F(n)\wedge F'(n+1)$ is a type $(n+1)$-space, hence $\Map(\Sigma^\infty F'(n+1),\Phi_{F(n)}(X))\simeq\Phi_{F(n)\wedge F'(n+1)}(X)\simeq0$. For the second fact, it suffices to show that $\Phi_{F(n)\wedge F'(n)}(X)\neq0$, which follows from \cref{Periodic_equiv_indep_type_n} since if the converse holds, then $X\to *$ would be a $v_n$-periodic equivalence, which implies $\Phi_{F(n)}(X)\simeq0$. 
	\end{proof}

	\begin{thm}\label{Pre_Bousfield_Kuhn_inv_tel_loc}
		For any pointed space $X$, for $0<m\leq n$ and any $d$, we have a canonical equivalence $\Phi_{F(m)}(X)\simeq\Phi_{F(m)}(L_{n,d}^fX)$
	\end{thm}
	\begin{proof}
		Recall that $L_{n,d}^f:\An_*\to L_n^f(\An_*)_{(p)}^{>d}$ is the composition $X\mapsto P_{F(n+1,d)}(\tau_{\geq d+1}X_{(p)})$. Hence it suffices to verify the three constructions do not change $\Phi_{F(m)}$. 

		Firstly, we prove $\tau_{\geq d+1}X\to X$ induces an equivalence $\Phi_{F(m)}(\tau_{\geq d+1}X)\simeq\Phi_{F(m)}(X)$, which is obvious since we have 
		\[ \Omega^\infty\Phi_{F(m)}(X)\simeq\ilim(\Hom_*(F(m),\Omega^{Nt}X)\to\Hom_*(F(m),\Omega^{Nt+t}X)\to\cdots) \]
		passes to a sufficiently large loopspace of $X$, which is not affected by the connective cover. 

		Then we prove that $X\to X_{(p)}$ induces an equivalence $\Phi_{F(m)}(X)\to\Phi_{F(m)}(X_{(p)})$. This is obtained by considering the pullback square of localizations $X\simeq X_{(p)}\times_{X_\mathbb{Q}}X[p^{-1}]$, which results in the pullback square
		% https://q.uiver.app/#q=WzAsNCxbMCwwLCJYIl0sWzEsMCwiWF97KHApfSJdLFswLDEsIlhbcF57LTF9XSJdLFsxLDEsIlhfXFxtYXRoYmJ7UX0iXSxbMCwxXSxbMSwzXSxbMCwyXSxbMiwzXSxbMCwzLCIiLDEseyJzdHlsZSI6eyJuYW1lIjoiY29ybmVyLWludmVyc2UifX1dXQ==
		\[\begin{tikzcd}[ampersand replacement=\&]
			{\Hom_*(F(m),X)} \& {\Hom_*(F(m),X_{(p)})} \\
			{\Hom_*(F(m),X[p^{-1}])} \& {\Hom_*(F(m),X_\mathbb{Q})}
			\arrow[from=1-1, to=1-2]
			\arrow[from=1-1, to=2-1]
			\arrow["\ulcorner"{anchor=center, pos=0.125}, draw=none, from=1-1, to=2-2]
			\arrow[from=1-2, to=2-2]
			\arrow[from=2-1, to=2-2]
		\end{tikzcd}\]

		Now by the universal property of localizations, we have 
		\[ \Hom_*(F(m),X[p^{-1}])\simeq\Hom_*(F(m)[p^{-1}],X[p^{-1}]) \] 
		and hence vanishes since $F(m)$ is of type $>0$. The likewise fact holds for $\Hom_*(F(m),X_\mathbb{Q})$, hence the bottom arrow is an equivalence, which proves that the top arrow is an equivalence. Thus we have $\Phi_{F(m)}(X)\simeq\Phi_{F(m)}(X_{(p)})$. 

		Finally it suffices to prove that for any $d$-connected space, the localization $X\to P_{F(n+1,d)}X$ induces an equivalence $\Phi_{F(m)}(X)\to\Phi_{F(m)}(P_{F(n+1,d)}(X))$. Again we may prove only that the $d$-connective cover of the zeroth space of the two spectra are equivalent. Firstly notice that $L_{T(m)}\Sigma^\infty F(n+1,d)\simeq0$ since $F(n+1)$ has type $n+1>m$. But by \cref{Pre_Bousfield_Kuhn_T_n_loc} we see that $\Phi_{F(m)}(X)$ is $T(m)$-local, hence we have 
		\[ \Hom_*(F(n+1,d),\Omega^\infty\Phi_{F(m)}(X))\simeq\Hom_*(\Sigma^\infty F(n+1,d),\Phi_{F(m)}(X))\simeq* \]
		and hence $\Omega^\infty\Phi_{F(m)}(X)$ is $F(n+1,d)$-null. Therefore the desired equivalence is an exchange map of $P_{F(n+1,d)}$ with the second factor of $\Hom$, hence we can make a few reductions as shown below (we remark here that $\tau_{\geq d+1}$ commutes with filtered colimits since it is a right adjoint). 
		\begin{align*}
			&\tau_{\geq d+1}\Omega^\infty\Phi_{F(m)}(X)\simeq\tau_{\geq d+1}\Omega^\infty\Phi_{F(m)}(P_{F(n+1,d)}X) \\
			\iff&\tau_{\geq d+1}P_{F(n+1,d)}\ilim(\Hom_*(\Sigma^{kt}F(m),X))\simeq\Omega^\infty\ilim(\Hom_*(\Sigma^{kt}F(m),P_{F(n+1,d)}X)) \\
			\iff&\ilim(\tau_{\geq d+1}(P_{F(n+1,d)}(\Hom_*(\Sigma^{kt}F(m),X))))\simeq\ilim(\tau_{\geq d+1}\Hom_*(\Sigma^{kt}F(m),P_{F(n+1,d)}X)) \\
			\iff&P_{F(n+1,d)}(\Hom_*(\Sigma^{kt}F(m),X))\simeq\Hom_*(\Sigma^{kt}F(m),P_{F(n+1,d)}X),\forall k
		\end{align*}
		But since $F(m)$ is finite, we may write $\Sigma^{kt}F(m)$ as a finite colimit of $S^0$. This makes \\
		$\Hom_*(\Sigma^{kt}F(m),X)$ a finite limit, which commutes with $P_{F(n+1,d)}$ by \cref{Tel_loc_fin_lim}. 
	\end{proof}

	\begin{prop}\label{Periodic_homotopy_grp_equiv_d_conn}
		Let $X$ be a $F(n+1,d)$-null pointed space. Then for any type $n$ space $F(n)$ equipped with a $v_n$-self map $v_n:\Sigma^tF(n)\to F(n)$, we have the map $\pi_k\Hom_*(F(n),X)\to v_n^{-1}\pi_k(X;F(n))$ an isomorphism for $k>d$. 
	\end{prop}
	\begin{proof}
		We claim that the transition maps in the colimit 
		\[ \Hom_*(F(n),X)\to\Hom_*(\Sigma^tF(n),X)\to\cdots \]
		are homotopy equivalences after passing to the $(d+1)$-connective covers, which proves the assertion directly by passing to the colimit. We claim the identity component of the homotopy fiber, that is, the identity component of $\Hom_*(\cof(v_n),X)$ is $(d-1)$-truncated. To prove this, note that 
		\[ \Omega^{d-1}\Hom_*(\cof(v_n),X)\simeq\Hom_*(\Sigma^{d-1}\cof(v_n),X)\simeq* \]
		where the last equivalence is by the fact that $X$ is $\Sigma^{d-1}\cof(v_n)$-null, since $\type\Sigma^{d-1}\cof(v_n)\geq n+1$, and $\conn\Sigma^{d-1}\cof(v_n)\geq d$, and by \cref{Bousfield_thm} we have proven the result. 
	\end{proof}

	\begin{thm}\label{Tel_loc_equiv_crit}
		Let $f:X\to Y$ be a map of pointed spaces. Then $L_{n,d}^f(f)$ is an equivalence iff for $0<m\leq n$, $f$ is a $v_m$-periodic equivalence and $(\pi_kX)_\mathbb{Q}\to(\pi_kY)_\mathbb{Q}$ is an equivalence for $k>d$. 
	\end{thm}
	\begin{proof}
		$\implies$: The first condition follows from \cref{Pre_Bousfield_Kuhn_inv_tel_loc}, and the second follows from the fact that $\fib(f)$ must be $F(n+1,d)$-acyclic, hence must be rationally trivial on degree greater than $d$ (the degrees below $d$ are being removed in the construction). 

		$\impliedby$: We apply induction on $n$. For the case where $n=0$, we have $L_0^f(X)\simeq(\tau_{\geq d+1}X)_\mathbb{Q}$, where the assumptions clearly imply $L_{n,d}^f(f)$ is an equivalence. For $n>0$, WLOG assume that $X$ and $Y$ are already $(d+1)$-connected, $p$-local and $F(n+1,d)$-null. Then by the assumption the fiber $F=\fib(f)$ is rationally trivial and $v_m$-periodic homotopy equivalent to a point for $0<m\leq n$. Let $F(n,d')$ be a type $n$ space of connectivity $d'$ equipped with a $v_n$-self map. By \cref{Periodic_homotopy_grp_equiv_d_conn}, the space $\Hom_*(F(n,d'),F)$ is $d$-truncated since it is $v_n$-homotopy equivalent to a point. By replacing $F(n,d')$ by a suitable suspension we may assume $\Hom_*(F(n,d'),F)\simeq*$. We moreover assume that $d'>d$. Therefore $F$ is $F(n,d')$-null. Then $\tau_{\geq d'+1}F$ is also $F(n,d')$-null since $\tau_{\geq d'+1}$ is a right adjoint, hence $\tau_{\geq d'+1}F$ can be regarded as an object in $L_{n-1}^f(\An_*)_{(p)}^{>d'}$. By the induction hypothesis on $n$, we conclude that $\tau_{\geq d'+1}F\simeq*$. Moreover, by assumption $\tau_{\geq d+1}F$ is rationally trivial and $F(n+1,d)$-null, by \cref{Tel_loc_conn_var_rational} we have $P_{F(n+1,d)}\tau_{\geq d'+1}F\simeq\tau_{\geq d+1}F$. Therefore we have $F$ is $d$-truncated. Moreover, since it is a fiber of a map of $d+1$-connective spaces, it must be $d$-connective. Therefore $F$ is concentrated at degree $d$, and hence $F\simeq K(G,d)$ for some $G$. Since $f$ is a rational equivalence on degree greater than $d$, $G$ must be torsion, and since $F$ is $p$-local, $G$ must be $p$-power torsion. Thus $K(G,d)$ is $F(n+1,d)$-periodic by \cref{Tel_loc_acyc_cond}. Thus $f$ is an equivalence since $X$ and $Y$ are $F(n+1,d)$-null. 
	\end{proof}

	\begin{coro}\label{Tel_loc_equiv_maps}
		The functor $L_{n,d}^f:\An_*\to L_n^f(\An_*)_{(p)}^{>d}$ exhibits the category $L_n^f(\An_*)_{(p)}^{>d}$ as the localization of $\An_*$ with respect to the maps satisfying the conditions in \cref{Tel_loc_equiv_crit}. 
	\end{coro}
	\begin{proof}
		Directly by the definition of localizations and \cref{Tel_loc_equiv_crit}. 
	\end{proof}

	\begin{lemma}\label{Pre_Bousfield_Kuhn_vanish_crit}
		For a $F(n)$-null space $X$, we have $\Phi_{F(m)}(X)\simeq0$ for all $m>n$. 
	\end{lemma}
	\begin{proof}
		By replacing $F(m)$ with a suspension, we may assume that the connectivity of $F(m)$ is greater than the connectivity of $F(n)$. Therefore by \cref{Bousfield_thm} we have $X$ is $F(m)$-null. Therefore the mapping space $\Hom_*(F(m),X)\simeq*$, and hence $\Phi_{F(m)}(X)\simeq0$. 
	\end{proof}

	\begin{defn}\label{Monochromatic_loc}
		We let $\An_*^{v_n}$ be the full subcategory of $L_n^f(\An_*)_{(p)}^{>d}$ consisting of $F(n,d)$-periodic objects. Explicitly, $\An_*^{v_n}$ is the category of spaces that are $(d+1)$-connective, $p$-local, $F(n+1,d)$-null, and $F(n,d)$-periodic. 

		We define a localization functor $M_n^f:\An_*\to\An_*^{v_n}$ by 
		\[ X\mapsto\tau_{\geq d+1}\fib(P_{F(n+1,d)}(\tau_{\geq d+1}X_{(p)})\to P_{F(n,d)}(\tau_{\geq d+1}X_{(p)})) \]
	\end{defn}

	\begin{coro}\label{Monochromatic_loc_v_n_equiv}
		A morphism in $\An_*^{v_n}$ is an equivalence iff it is an $v_n$-periodic homotopy equivalence. 
	\end{coro}
	\begin{proof}
		By \cref{Tel_loc_equiv_crit}. 
	\end{proof}

	\begin{coro}\label{Monochromatic_loc_equiv_maps}
		The functor $M_n^f:\An_*\to\An_*^{v_n}$ exhibits $\An_*^{v_n}$ as the localization of $\An_*$ with respect to the $v_n$-periodic equivalences. 
	\end{coro}
	\begin{proof}
		Directly by \cref{Monochromatic_loc_v_n_equiv}
	\end{proof}

	\begin{coro}\label{Tel_susp_loop}
		The stabilization $\Sp(\An_*^{v_n})$ can be identified with $\Sp_{T(n)}$. We denote the associated suspension spectrum and infinite loopspace functor by $\Sigma^\infty_{T(n)}$ and $\Omega^\infty_{T(n)}$. Moreover, we have $\Sigma^\infty_{T(n)}\simeq L_{T(n)}\Sigma^\infty$ (under a certain inclusion $\An_*^{v_n}\inj\An_*$) and $\Omega^\infty_{T(n)}\simeq\Omega^\infty\circ\tau_{\geq d_{n+1}+1}$, where $d_{n+1}$ depends on the inclusion chosen $\An_*^{v_n}\inj\An_*$. 
	\end{coro}
	\begin{proof}
		By verifying universal properties, $\Sp(\mathcal{C}[W^{-1}])$ can be identified with the category \\
		$\Sp(\mathcal{C})[\Sigma^\infty W^{-1}]$. Therefore $\Sp(\An_*^{v_n})$ is the category obtained by inverting the $v_n$-periodic equivalences in $\Sp$, which is obviously $\Sp_{T(n)}$. 
	\end{proof}

	The above corollary shows that the nature of the category $\An_*^{v_n}$ do not depend on the value of $d$ in $L_n^f(\An_*)_{(p)}^{>d}$, since $\An_*^{v_n}$ is uniquely determined by its universal property as a localization. However, what $d$ does affect is the inclusion functor $\An_*^{v_n}\inj\An_*$, which must have image in $(d+1)$-connected spaces. 

	\begin{thm}\label{Pre_Bousfield_Kuhn_inv_mono_loc}
		Let $F(n)$ be a type $n$ space equipped with a $v_n$-self map. Then the natural transformation $X\to M_n^fX$ induces an equivalence $\Phi_{F(n)}(X)\to\Phi_{F(n)}(M_n^fX)$. 
	\end{thm}
	\begin{proof}
		Recall in \cref{Pre_Bousfield_Kuhn_inv_tel_loc}, the natural transformations $X\to P_AX$, $X\to X_{(p)}$, and $\tau_{\geq d+1}X\to X$ are send to equivalences by $\Phi_{F(n)}$. Therefore it suffices to show that for a $F(n+1,d)$-null space $X$, the construction $X\mapsto\fib(X\to P_{F(n,d)}X)$ does not change the $v_n$-periodic homotopy groups of $X$, which is obvious since $P_{F(n,d)}X$ has trivial $v_n$-periodic homotopy groups. 
	\end{proof}

\section{The construction of the Bousfield Kuhn functor}

	In the last chapter we introduced the notion $\Phi_{F(n)}(X)$, which is exactly the information of the $v_n$-periodic homotopy groups of $X$. However, this construction (and also the $v_n$-periodic homotopy groups) depend on the choice of the type $n$ $p$-local space $F(n)$. In this section, we introduce the notion of the Bousfield-Kuhn functor, which unifies the construction $\Phi_{F(n)}(X)$. Namely, we prove the following theorem. 
	
	\begin{thm}\label{Bousfield_Kuhn_fun}
		There is a unique functor $\Phi:\An_*\to\Sp_{T(n)}$ such that for any finite $p$-local type $n$ space $F(n)$, we have $\Phi_{F(n)}(X)\simeq\Map(F(n),\Phi(X))$. 
	\end{thm}
	\begin{proof}
		We firstly notice that the notion $\Phi_{F(n)}(X)$ can be extended naturally to a functor $\Phi_-(-):(\Sp_n)^{\op}\to\Fun(\An_*,\Sp_{T(n)})$, where $\Sp_n$ is the subcategory of $p$-local finite spectra of type $\geq n$. This is essentially by \cref{Pre_Bousfield_Kuhn_susp_power} and \cref{Periodicity_thm_lem_5}. Then we consider $F:(\Sp_\mathrm{fin})^{\op}\to\Fun(\An_*,\Sp_{T(n)})$ as the right Kan extension of $\Phi_-(-)$. Since $\Sp_{T(n)}$ is closed under limits, the right Kan extension is well defined. Then we let $\Phi=F(\mathbb{S}_{(p)})$, namely by the construction of the right Kan extension, 
		\[ \Phi(X)=\plim\Phi_E(X) \]
		where the limit is taken on the over category $(\Sp_n^{\op})_{\mathbb{S}_{(p)}/}$. Then by \cref{Pre_Bousfield_Kuhn_cofib}, we have $\Phi_-(-):(\Sp_n)^{\op}\to\Fun(\An_*\to\Sp_{T(n)})$ is left exact. Therefore $F$ is also left exact, and since 
		\[ F(\mathbb{S}_{(p)})(X)=\Phi(X)\simeq\Map(\mathbb{S}_{(p)},\Phi(X)) \]

		By the definition of the Kan extension and the definition of extension by limits we have 
		\[ F(F(n))(X)=\Phi_{F(n)}(X)=\Map(F(n),\Phi(X)) \]

		Then we prove uniqueness. Suppose $\Phi':\An_*\to\Sp_{T(n)}$ is another functor satisfying the given conditions, then let $F':(\Sp_{\mathrm{fin}})^{\op}\to\Fun(\An_*,\Sp_{T(n)})$ by letting $F'(E)(X)=\Map(E,\Phi'(X))$. Then by the universal property of the Kan extension we obtain a natural transformation $F'\to F$, which is an equivalence on the full subcategory spanned by spectra of type $\geq n$. Therefore for some type $n$ spectrum $F(n)$, we have $\Map(F(n),\Phi(X))\simeq\Map(F(n),\Phi'(X))$. By taking duals we have $\dual F(n)\otimes\Phi(X)\simeq \dual F(n)\otimes\Phi'(X)$. By passing to the colimit we have $\Phi(X)\to\Phi'(X)$ is a $T(n)$-equivalence. But since both $\Phi(X)$ and $\Phi'(X)$ are $T(n)$-local, they must be equivalent. 
	\end{proof}

	\begin{remark}
		In practice, we often find a cofinal system of type $n$ spectra 
		\[ E_0\to E_1\to\cdots \]
		among all type $n$ spectra with a map to $\mathbb{S}_{(p)}$, and use $\plim_k\Phi_{E_k}(X)$ to calculate $\Phi(X)$. 
	\end{remark}

	\begin{prop}\label{Bousfield_Kuhn_lex}
		The functor $\Phi:\An_*\to\Sp_{T(n)}$ is left exact. 
	\end{prop}
	\begin{proof}
		It suffices to prove $\Phi_{F(n)}:\An_*\to\Sp_{T(n)}$ is left exact, where $F(n)$ is some type $n$ space with a $v_n$-self map. Since limits of spectra are computed spacewise, it suffices to prove $\Omega^{\infty+k}\Phi_{F(n)}$ is left exact. Then since $\Phi_{F(n)}$ is a periodic spectrum, it suffices to show $\Omega^\infty\Phi_{F(n)}$ is left exact. But by construction, 
		\[ \Omega^\infty\Phi_{F(n)}X=\ilim(\Hom_*(F(n),X)\to\Hom_*(\Sigma^tF(n),X)\to\cdots) \]
		which is clearly right exact since finite limits commutes with filtered colimits. 
	\end{proof}

	\begin{prop}\label{Bousfield_Kuhn_v_n_equiv}
		A map $f:X\to Y$ of pointed spaces is a $v_n$-periodic homotopy equivalence iff $\Phi(f):\Phi(X)\to\Phi(Y)$ is an equivalence. 
	\end{prop}
	\begin{proof}
		$\implies$ is obvious since by definition $v_n$-periodic homotopy equivalences are the maps send by $\Phi_{F(n)}$ to equivalences. Therefore by passing to limits we have $\Phi(f)$ an equivalence. 

		$\impliedby$ is also obvious since $\Phi_{F(n)}(X)\simeq\Map(F(n),\Phi(X))$. 
	\end{proof}
	
	\begin{coro}\label{Bousfield_Kuhn_source_v_n}
		The functor $\Phi:\An_*\to\Sp_{T(n)}$ factors through the localization $\An_*\to\An_*^{v_n}$. 
	\end{coro}
	\begin{proof}
		Directly by \cref{Monochromatic_loc_v_n_equiv} and \cref{Bousfield_Kuhn_v_n_equiv}. 
	\end{proof}

	\begin{thm}\label{Bousfield_Kuhn_adjoint}
		The functor $\Phi:\An_*^{v_n}\to\Sp_{T(n)}$ admits a left adjoint $\Theta:\Sp_{T(n)}\to\An_*^{v_n}$. 
	\end{thm}
	\begin{proof}
		By the construction of $\Phi$ as a right Kan extension, it suffices to prove $\Phi_{F(n)}$ for $F(n)$ a type $n$ space with a $v_n$-self map $v_n:\Sigma^tF(n)\to F(n)$ has a left adjoint $\Theta_{F(n)}$ (then we may construct $\Theta$ as a colimit of $\Theta_{F(n)}$). Explicitly, we show that for each spectrum $Z$, the functor $X\mapsto\Hom(Z,\Phi_{F(n)}(X))$ is corepresented by some $\Theta_{F(n)}(Z)\in\An_*^{v_n}$. 

		Since left adjoints preserves colimits, it suffices to prove $\Theta_{F(n)}(Z)$ exists for $Z=\mathbb{S}^{kt}$ with $k$ any integer (possibly negative). Using $t$-periodicity of $\Phi_{F(n)}(X)$ it suffices to prove the case for $Z=\mathbb{S}$, that is, the functor $X\mapsto\Omega^\infty\Phi_{F(n)}(X)$ is corepresentable. But by \cref{Periodic_homotopy_grp_equiv_d_conn} we have the sequence 
		\[ \Hom_*(F(n),X)\to\Hom_*(\Sigma^tF(n),X)\to\cdots \]
		is eventually constant since it is constant after taking a $d$-connective cover for a $F(n+1,d)$-null space $X$. Thus we only need to show that the functor $X\mapsto\Hom_*(\Sigma^{ct}F(n),X)$ is corepresentable for some $c\gg0$ (since $X\in\An_*^{v_n}$, which is also why the theorem does not hold if we regard $\Phi$ as a functor $\An_*\to\Sp$). 

		We claim this is corepresented by $P_{F(n+1,d)}(\Sigma^{ct}F(n))$. This space obviously lies in \\
		$L_n^f(\An_*)_{(p)}^{>d}$, and for $c\gg0$, $P_{F(n,d)}(\Sigma^{ct}F(n))\simeq*$ by \cref{Bousfield_thm} and the definition of the unstable Bousfield class. Hence $P_{F(n+1,d)}(\Sigma^{ct}F(n))\in\An_*^{v_n}$, and represents the functor, which proves the theorem. 
	\end{proof}

	The most fundamental property of the Bousfield Kuhn functor is that the $T(n)$-localization functor of spectra factors through it. 
	
	\begin{thm}\label{Bousfield_Kuhn_tel_loc_fac}
		Let $X$ be a spectrum. Then we have a canonical equivalence $\Phi\Omega^\infty(X)\simeq L_{T(n)}(X)$. 
	\end{thm}
	\begin{proof}
		We first calculate $\Phi_{F(n)}\Omega^\infty(X)$ where $F(n)$ is a type $n$ space equipped with a $v_n$-self map. We have 
		\begin{align*}
			&\Omega^\infty\Phi_{F(n)}(\Omega^\infty X) \\
			\simeq&\ilim(\Hom_*(F(n),\Omega^\infty X)\to\Hom_*(\Sigma^tF(n),\Omega^\infty X)\to\cdots) \\
			\simeq&\ilim(\Hom(\Sigma^\infty F(n),X)\to\Hom(\Sigma^{\infty+t}F(n),X)\to\cdots) \\
			\simeq&\ilim(\Omega^\infty D\Sigma^\infty F(n)\otimes X\to\Omega^{\infty-t}D\Sigma^\infty F(n)\otimes X\to\cdots) \\
			\simeq&\Omega^\infty \dual F(n)[v_n^{-1}]\otimes X \\
			\simeq&\Omega^\infty \dual F(n)[v_n^{-1}]\otimes L_{T(n)}X 
		\end{align*}
		where the fourth equivalence is since $\Omega^\infty$, as a right adjoint, commutes with filtered colimits, and the final equivalence is since $\dual F(n)[v_n^{-1}]$ is a telescope of type $n$, hence $\dual F(n)[v_n^{-1}]$ calculates the $T(n)$-homology of $X$ for some $X$, hence it is safe to replace $X$ by $L_{T(n)}(X)$. However, when $X$ is itself $T(n)$-local, the sequence 
		\[ \Hom(\Sigma^\infty F(n),X)\to\Hom(\Sigma^{\infty+t}F(n),X)\to\cdots \]
		is actually constant since the fiber of the map is $\Hom(\Sigma^\infty F(n)/v_n,X)$, and the previous spectrum is a spectrum of type $n+1$, and since $X$ is $T(n)$-local hence $F(n+1)$-null, this fiber is null. Therefore we have 
		\[ \Omega^\infty\Phi_{F(n)}(\Omega^\infty X)\simeq\Omega^\infty\Map(\Sigma^\infty F(n),L_{T(n)}X) \]
		and since both sides are periodic, we have $\Phi_{F(n)}(\Omega^\infty X)\simeq\Map(\Sigma^\infty F(n),L_{T(n)}X)$. This clearly holds after replacing $\Sigma^\infty F(n)$ by any type $n$ spectrum. 

		Thus the conclusion follows from 
		\[ \Phi(\Omega^\infty X)\simeq\plim_E\Phi_E(\Omega^\infty X)\simeq\plim_E\Map(E,L_{T(n)}X)\simeq L_{T(n)}X \]
	\end{proof}

	\begin{coro}\label{Bousfield_Kuhn_susp_loop_adj}
		We have two pairs of adjunctions 
		% https://q.uiver.app/#q=WzAsMyxbMiwwLCJcXFNwX3tUKG4pfSJdLFsxLDAsIlxcQW5fKl57dl9ufSJdLFswLDAsIlxcU3Bfe1Qobil9Il0sWzAsMSwiXFxPbWVnYV5cXGluZnR5X3tUKG4pfSIsMCx7Im9mZnNldCI6LTF9XSxbMSwyLCJcXFBoaSIsMCx7Im9mZnNldCI6LTF9XSxbMiwxLCJcXFRoZXRhIiwwLHsib2Zmc2V0IjotMX1dLFsxLDAsIlxcU2lnbWFeXFxpbmZ0eV97VChuKX0iLDAseyJvZmZzZXQiOi0xfV1d
		\[\begin{tikzcd}[ampersand replacement=\&]
			{\Sp_{T(n)}} \& {\An_*^{v_n}} \& {\Sp_{T(n)}}
			\arrow["\Theta", shift left, from=1-1, to=1-2]
			\arrow["\Phi", shift left, from=1-2, to=1-1]
			\arrow["{\Sigma^\infty_{T(n)}}", shift left, from=1-2, to=1-3]
			\arrow["{\Omega^\infty_{T(n)}}", shift left, from=1-3, to=1-2]
		\end{tikzcd}\]
		such that the horizontal compositions are the identities. 
	\end{coro}
	\begin{proof}
		\cref{Bousfield_Kuhn_tel_loc_fac} implies that the bottom composition is the identity, and the top row is left adjoint to the bottom row, hence the identity. 
	\end{proof}

\section{Monadicity of the Bousfield Kuhn functor}

	We first recall a few facts on monads. 

	\begin{defn}\label{Infty_monads}
		Recall that traditionally (or in the 1-categorical sense), a monad on a 1-category $\mathcal{C}$ is an endofunctor $T:\mathcal{C}\to\mathcal{C}$ equipped with a natural transformation $\mu:T\circ T\to T$ and a natural transformtion $\eta:\id_\mathcal{C}\to T$ such that the following diagrams commute. 
		% https://q.uiver.app/#q=WzAsOCxbMCwwLCJUXFxjaXJjIFRcXGNpcmMgVCJdLFsxLDAsIlRcXGNpcmMgVCJdLFswLDEsIlRcXGNpcmMgVCJdLFsxLDEsIlQiXSxbMiwwLCJUIl0sWzMsMCwiVFxcY2lyYyBUIl0sWzQsMCwiVCJdLFszLDEsIlQiXSxbMCwxLCJcXG11XFxjaXJjIDFfVCJdLFsxLDMsIlxcbXUiXSxbMCwyLCIxX1RcXGNpcmMgXFxtdSIsMl0sWzIsMywiXFxtdSIsMl0sWzQsNSwiMV9UXFxjaXJjXFxldGEiXSxbNSw3LCJcXG11IiwyXSxbNiw1LCJcXGV0YVxcY2lyYyAxX1QiLDJdLFs2LDcsIjFfVCJdLFs0LDcsIjFfVCIsMl1d
		\[\begin{tikzcd}[ampersand replacement=\&]
			{T\circ T\circ T} \& {T\circ T} \& T \& {T\circ T} \& T \\
			{T\circ T} \& T \&\& T
			\arrow["{\mu\circ 1_T}", from=1-1, to=1-2]
			\arrow["{1_T\circ \mu}"', from=1-1, to=2-1]
			\arrow["\mu", from=1-2, to=2-2]
			\arrow["{1_T\circ\eta}", from=1-3, to=1-4]
			\arrow["{1_T}"', from=1-3, to=2-4]
			\arrow["\mu"', from=1-4, to=2-4]
			\arrow["{\eta\circ 1_T}"', from=1-5, to=1-4]
			\arrow["{1_T}", from=1-5, to=2-4]
			\arrow["\mu"', from=2-1, to=2-2]
		\end{tikzcd}\]
		where $1_T$ denotes the trivial natural transformation $1_T:T\to T$ given by $\id$. Finally, an algebra of a monad is an object of $\mathcal{C}$, say $X$, equipped with a morphism $m:TX\to X$ such that the following diagrams commute 
		% https://q.uiver.app/#q=WzAsNyxbMCwwLCJUXFxjaXJjIFRYIl0sWzEsMCwiVFgiXSxbMCwxLCJUWCJdLFsxLDEsIlgiXSxbMiwwLCJYIl0sWzMsMCwiVFgiXSxbMywxLCJYIl0sWzAsMSwiXFxtdVxcY2lyYyAxX1giXSxbMSwzLCJtIl0sWzAsMiwiMV9UXFxjaXJjIG0iLDJdLFsyLDMsIm0iLDJdLFs0LDUsIlxcZXRhX1giXSxbNSw2LCJtIl0sWzQsNiwiMV9YIiwyXV0=
		\[\begin{tikzcd}[ampersand replacement=\&]
			{T\circ TX} \& TX \& X \& TX \\
			TX \& X \&\& X
			\arrow["{\mu\circ 1_X}", from=1-1, to=1-2]
			\arrow["{1_T\circ m}"', from=1-1, to=2-1]
			\arrow["m", from=1-2, to=2-2]
			\arrow["{\eta_X}", from=1-3, to=1-4]
			\arrow["{1_X}"', from=1-3, to=2-4]
			\arrow["m", from=1-4, to=2-4]
			\arrow["m"', from=2-1, to=2-2]
		\end{tikzcd}\]
		where $1_X$ is the identity map on $X$. 

		The constructions immediately generalizes to $\infty$-categories. Note that the above definitions abstracts to the following. Let $\mathcal{C}$ be an $\infty$-category, then $\Fun(\mathcal{C},\mathcal{C})\simeq\Hom_{\Cat}(\mathcal{C},\mathcal{C})$ has a tautological structure of an $\mathbb{E}_1$ monoidal category with tensor product given by composition of functors (here $\Cat$ is regarded as a $(\infty,2)$-category). Then $\mathcal{C}$ is a category left tensored over $\Fun(\mathcal{C},\mathcal{C})$ given by applying an endofunctor on an object. Therefore, the above definition of monads generalizes as follows. We define a monad $T$ on a category $\mathcal{C}$ to be an $\mathbb{E}_1$-algebra in $\Fun(\mathcal{C},\mathcal{C})$, and a $T$-algebra $X$ to be an object of $\LMod_T(\mathcal{C})$. 
	\end{defn}

	\begin{const}\label{Adj_EM_K_fun_defn}
		For a monad $T$ on $\mathcal{C}$, notice that we have the following adjunction 
		\[ T:\mathcal{C}\adj\LMod_T(\mathcal{C}):U \]
		which also recovers the monad by the composition $UT$.  
		For a pair of adjunction $F:\mathcal{C}\adj\mathcal{D}:G$, we have $GF$ a monad on $\mathcal{C}$, with its category of algebras equivalent to the essential image of $G$. 

		Explicitly, the $\mathbb{E}_1$-algebra structure of $GF$ has the multiplication map $1_G\circ\varepsilon\circ 1_F:GFGF\to GF$ and the unit map $\eta:\id_\mathcal{C}\to GF$. Furthermore, notice that for any element $Y\in\mathcal{D}$, $GY$ is a $GF$-algebra with the action given by $1_G\circ\varepsilon\circ 1_Y:GFGY\to GY$. Hence $G$ upgrades to a comparison functor (the Eilenberg-Moore functor) $E:\mathcal{D}\to\LMod_{GF}(\mathcal{C})$, and we say $GF$ is monadic iff the previous comparison functor is an equivalence. 

		We have the further definition of the Kleisli category, or the category of free $GF$-modules $\LMod_{GF}^{\mathrm{free}}(\mathcal{C})$, defined to be the full subcategory of $\LMod_{GF}(\mathcal{C})$ consisting of the $GF$-modules of the form $GFX$ equipped with the natural module structure. By the definition of adjunctions, a morphism of $GF$-modules $GFX\to GFY$ is equivalent to the data of a morphism $X\to GFY$ as objects of $\mathcal{C}$. Therefore, the Kleisli category is equivalent to a category having the same objects as $\mathcal{C}$, and a morphism $X\to Y$ is defined to be a morphism $X\to GFY$ in $\mathcal{C}$ (this reduces our construction to the usual construction of Kleisli categories). This allows us to define a map 
		\[ K:\LMod_{GF}^{\mathrm{free}}(\mathcal{C})\to\mathcal{D} \] 
		given by sending an object $X\in\LMod_{GF}^{\mathrm{free}}(\mathcal{C})$ via the second construction as an object in $\mathcal{C}$ to $FX$, and a morphism $X\to Y$ regarded as a morphism $X\to GFY$ is send to $FX\to FY$ given by the defintion of adjunctions. By this identification we have obviously that $K$ is fully faithful. 

		In the aobve constructions, we see that we have $E\circ K:\LMod_{GF}^{\mathrm{free}}(\mathcal{C})\to\LMod_{GF}(\mathcal{C})$ is naturally equivalent to the inclusion functor of the free modules into all modules. In particular, $E\circ K$ is fully faithful. 
	\end{const}

	We have the following fundamental theorem of Barr-Beck on the monadicity of adjoint functors, generalized by Lurie to the $\infty$-categorical case, whose proof we will omit. 

	\begin{thm}[Barr-Beck-Lurie]\label{Barr_Beck_Lurie}
		The adjunction $F:\mathcal{C}\adj\mathcal{D}:G$ is monadic iff the two following conditions holds. 
		\begin{enumerate}
			\item $G$ is conservative. 
			\item For a simplicial object $V$ of $\mathcal{D}$, if it is send to a split simplicial object in $\mathcal{C}$ by $G$, then it admit a colimit in $\mathcal{D}$, and this colimit is preserved by $G$. 
		\end{enumerate}
		In particular, if $\mathcal{D}$ admits all geometric realizations, all of which is preserved by $G$, then the second condition clearly holds. 
	\end{thm}

	We can now state the first main result of this review. 

	\begin{thm}[Eldred-Heuts-Mathew-Meier]\label{Monadicity_Bousfield_Kuhn}
		The adjunction $\Theta:\Sp_{T(n)}\adj\An_*^{v_n}:\Phi$ is monadic. 
	\end{thm}

	Before this, we need some lemmas on geometric realizations. 

	\begin{lemma}\label{Bousfield_Friedlander_crit}
		For a diagram of simplicial spaces $X_\bullet\to B_\bullet\from Y_\bullet$ with $B_n$ is connected for all $n$, then the natural map $|X_\bullet\times_{B_\bullet}Y_\bullet|\to|X_\bullet|\times_{|B_\bullet|}|Y_\bullet|$ is an equivalence. 
	\end{lemma}
	\begin{proof}
		This is a famous result called the Bousfield-Friedlander criterion in simplicial homotopy theory, hence we will omit the proof here. \todo{Reference!}
	\end{proof}

	\begin{lemma}\label{Hom_geom_realization}
		Let $V\in\An_*$ be a finite $d$-dimensional CW complex. Then the functor $\Hom_*(V,-):\An_*^{\geq d}\to\An_*$ preserves geometric realizations. 
	\end{lemma}
	\begin{proof}
		Consider a simplicial space $X_\bullet$ with each $X_n$ $d$-connective. Then we consider the collection $T$ of all finite spaces $V$ such that $|\Hom_*(V,X_\bullet)|\simeq\Hom_*(V,|X_\bullet|)$. Note that the category $\An_*$ is a topos, and hence $T$ is closed under coproducts (wedge sums). 

		We apply induction to show all spaces of dimension $\leq d$ is in $T$. Firstly by definition $S^0\in T$, and since $T$ is closed under wedge sums, $0$-dimensional spaces are all in $T$. Assume now that for $n\leq d-1$ that all $n$-dimensional spaces are in $T$, then for any space $Y$ of dimension $n+1$, we have a cofiber sequence $\bigvee S^n\to\tau_{\leq n}Y\to Y$. Moreover, since $n\leq d-1$, $\Hom_*(S^n,X_\bullet)$ is connected, and hence by \cref{Bousfield_Friedlander_crit} and that $\bigvee S^n$ and $\tau_{\leq n}Y$ are in $T$, we have $Y\in T$. 
	\end{proof}

	\begin{proof}
		We now apply \cref{Barr_Beck_Lurie} to prove \cref{Monadicity_Bousfield_Kuhn}. Firstly $\Phi$ is conservative by \cref{Bousfield_Kuhn_v_n_equiv} and \cref{Monochromatic_loc_v_n_equiv}. Hence it suffices to prove that $\Phi$ preserves geometric realizations. 

		First note that the colimit in $\Sp_{T(n)}$ is obtained by first taking the colimit in $\Sp$ then take the $T(n)$-localization. This allows us to regard $\Phi$ as a functor $\An_*^{v_n}\to\Sp$. Let $X_\bullet$ be a simplicial object in $\An_*^{v_n}$, we prove $|\Phi(X_\bullet)|\to\Phi(|X_\bullet|)$ is a $T(n)$-equivalence. Since $T(n)$ can be obtained by inverting the $v_n$-self map of any type $n$ spectrum, and the dual of a type $n$ spectrum is still type $n$, it suffices to prove 
		\[ \Map(F(n),|\Phi(X_\bullet)|)\to\Map(F(n),\Phi(|X_\bullet|)) \] 
		is an equivalence. Now since in the category of spectra fiber sequences coincide with the cofiber sequences, the fact that geometric realization is a colimit and hence preserves every colimit implies that it commutes with finite limits. Now $\Map(F(n),-)$ is a finite limit since $F(n)$ is finite, hence the left hand side can be identified with $|\Map(F(n),\Phi(X_\bullet))|$. Therefore it suffices to prove 
		\[ |\Map(F(n),\Phi(X_\bullet))|to\Map(F(n),\Phi(|X_\bullet|)) \]
		is an equivalence. By \cref{Bousfield_Kuhn_fun} and choosing $F(n)$ to be of the form $\Sigma^\infty F(n)$ for a type $n$ space $F(n)$, it suffices to prove 
		\[ |\Phi_{F(n)}(X_\bullet)|\to\Phi_{F(n)}(|X_\bullet|) \]
		is an equivalence. Then, notice that the subcategory $\An_*^{v_n}\subseteq L_n^f(\An_*)_{(p)}^{>d}$ is closed under colimits, and the colimits in $L_n^f(\An_*)_{(p)}^{>d}$ are computed by first taking the colimits in $\An_*$ and then take the localization $L_{n,d}^f$. Finally we notice that the map $L_{n,d}^f:(\An_*)_{(p)}^{>d}\to L_n^f(\An_*)_{(p)}^{>d}$ induces an equivalence on $v_n$-periodic homotopy, hence induces an equivalence after applying $\Phi_{F(n)}$. Therefore it suffices to show that $\Phi_{F(n)}:(\An_*)_{(p)}^{>d}\to\Sp$ preserves geometric realizations for some chosen $d$. 

		Finally we expand the definition of $\Phi_{F(n)}(X_\bullet)$ for $X_n\in(\An_*)_{(p)}^{>d}$. Suppose that the $v_n$-self map of $F(n)$ is $v_n:\Sigma^tF(n)\to F(n)$, we notice that 
		\[ \Phi_{F(n)}(X_\bullet)\simeq\ilim(\tau_{\geq0}\Phi_{F(n)}(X_\bullet)\to\Omega^t\tau_{\geq0}\Phi_{F(n)}(X_\bullet)\to\cdots) \]
		hence it suffices to see that $\tau_{\geq0}\Phi_{F(n)}$ preserves geometric realizations. But note that the category of $\mathbb{E}_\infty$-groups are closed under filtered colimits as the subcategory of $\An_*$, hence under the identification of the category of connective spectra and $\mathbb{E}_\infty$-groups, and the fact that 
		\[ \tau_{\geq0}\Phi_{F(n)}(X_\bullet)\simeq B^\infty\ilim(\Hom_*(F(n),X_\bullet)\to\Hom_*(\Sigma^tF(n),X_\bullet)\to\cdots) \]
		it suffices to show 
		\[ |\Hom_*(F(n),X_\bullet)|to\Hom_*(F(n),|X_\bullet|) \]
		is an isomorphism. But since the category $\An_*^{v_n}$ is independent of the choice of $d$ in $L_n^f(\An_*)_{(p)}^{>d}$, we may choose $d\gg0$ such that $d>\dim F(n)$, and hence the theorem follows from \cref{Hom_geom_realization}
	\end{proof}

	As a direct corollary, $\Theta$ computes the free algebra (of the monad $\Phi\Theta$) of a $T(n)$-local spectra. Therefore the computation of the functor $\Theta$ may appear as an important question. In fact, for a $T(n)$-local finite spectrum, we have the following way of computing $\Theta$. 

	\begin{prop}\label{Val_theta_fin_sp}
		Let $F(n)$ be a $(d-2)$-connected finite space of type $n$ with a $v_n$-self map $v:\Sigma^tF(n)\to F(n)$, then we have a canonical equivalence 
		\[ \Theta(L_{T(n)}\Sigma^{\infty+2}F(n))\simeq L_{n,d}^f\Sigma^2F(n) \]
	\end{prop}
	\begin{proof}
		Consider the space $F(n)/v$. Then by \cref{Bousfield_thm}, $\Sigma^iF(n)/v$ a suspension space having type $n+1$ and connectivity $d+i-1$, hence it is $F(n,d)$-periodic. Thus
		\[ \Hom_*(L_{n,d}^f\Sigma^iF(n)/v,X)\simeq\Hom_*(\Sigma^iF(n)/v,X)\simeq* \]

		Then consider the cofiber sequence 
		\[ \Sigma F(n)/v\to\Sigma^{t+2}F(n)\to\Sigma^tF(n)\to\Sigma^2F(n)/v \]

		Let $X\in\An_*^{v_n}$. By the above, we have 
		\[ \Hom_*(\Sigma^2F(n),X)\simeq\Hom_*(\Sigma^{t+2}F(n),X)\simeq\cdots \]

		Thus for any $X\in\An_*^{v_n}$, we have 
		\begin{align*}
			&\Hom_*(\Theta(L_{T(n)}\Sigma^{\infty+2}F(n)),X) \\
			\simeq&\Hom(L_{T(n)}\Sigma^{\infty+2}F(n),\Phi(X)) \\
			\simeq&\Omega^\infty\Phi_{\Sigma^2F(n)}(X) \\
			\simeq&\Hom_*(\Sigma^2F(n),X) \\
			\simeq&\Hom_*(L_{n,d}^f\Sigma^2F(n),X)
		\end{align*}
	\end{proof}

\section{Calculus and dual calculus of functors}

	Recall that we have the general philosophy that there is a monoidal (not generally fully faithful) embedding $\Fun(\Core\Fin,\mathcal{C})\to\Fun(\mathcal{C},\mathcal{C})$ given by 
	\[ C\mapsto(X\mapsto\bigoplus(C(n)\otimes X^{\otimes n})_{hS_n}) \] 
	for any category $\mathcal{C}$. Thus given an operad in the classical sense, that is, a symmetric sequence $C\in\Fun(\Core\Fin,\mathcal{C})$ which is a $\mathbb{E}_1$-monoid under the composition product, it gives rise to a monad. The essense of (dual) Goodwillie calculus which is used here finds, in some sense, adjoints to this construction. In this section, $\mathcal{C}$ and $\mathcal{D}$ are \emph{pointed} (or even stable) categories good enough so that they have all limits and colimits that are possibly implicitly needed in the conditions. 

	\begin{defn}\label{Tot_fib}
		For an $S$-cube $X:P(S)\to\mathcal{C}$, we define its total fiber 
		\[ \tfib X=\fib(X(\varnothing)\to\plim_{T\subseteq S,T\neq\varnothing}X(T)) \]
		and dually its total cofiber 
		\[ \tcof X=\cof(\ilim_{T\subseteq S,T\neq S}X(T)\to X(S)) \]
	\end{defn}

	\begin{coro}\label{Tot_fib_iterative}
		The construction of total (co)fibers are iterative. That is, given a cube $X:P(S)\to\mathcal{C}$ and $s\in S$, we have $\tfib(X)$ identified with the total fiber of the following cube 
		\[ Y:P(S-\{s\})\to\mathcal{C}, Y(T)=\fib(Y(T)\to Y(T\cup\{s\})) \]

		The dual statement holds for the total cofiber. 
	\end{coro}
	\begin{proof}
		By directly verifying the universal property. 
	\end{proof}

	\begin{defn}\label{Exc_fun}
		Consider a finite set $S$, and let $P(S)$ be the poset of subsetes of $S$ ordered by inclusion. We define $P_{\leq i}(S)$ to be the subset of $P(S)$ consisting of subsets having cardinality smaller or equal than $i$ and $P_{>i}(S)$ to be the subset of $P(S)$ consisting of subsets having cardinality greater than $i$. 

		We define an $S$-cube $F:P(S)\to\mathcal{C}$ (where $S$ is a finite set) to be strongly coCartesian if it is the left Kan extension of the subdiagram $F:P_{\leq 1}(S)\to\mathcal{C}$. We define a cube to be Cartesian if it is the right Kan extension of the subdiagram $F:P_{>0}(S)\to\mathcal{C}$. 

		We define a functor $F:\mathcal{C}\to\mathcal{D}$ to be $n$-excisive if it takes strongly coCartesian $[n]$-cubes to Cartesian $[n]$-cubes (here $[n]=\{0,1,\cdots,n\}$). Clearly we have a functor is $n$-excisive implies that it is $n+1$-excisive. We define the full subcategory $\Exc_n(\mathcal{C},\mathcal{D})$ to be the full subcategory of $\Fun(\mathcal{C},\mathcal{D})$ spanned by the $n$-excisive functors. 

		Also, we form the dual of the above construction. We have the dual definition of strongly Cartesian cubes and coCartesian cubes. We define a functor to be $n$-coexcisive if it takes strongly Cartesian $[n]$-cubes to coCartesian $[n]$-cubes. Dually we also have a functor is $n$-coexcisive implies that it is $n+1$-excisive. We define the category of coexcisive functors to be $\Exc^n(\mathcal{C},\mathcal{D})$. 

%		We have left adjoints and right adjoints to the inclusion of the excisive and coexcisive functors into the category of all functors, which we denote by $P_k:\Fun(\mathcal{C},\mathcal{D})\to\Exc_n(\mathcal{C},\mathcal{D})$ and $P^k:\Fun(\mathcal{C},\mathcal{D})\to\Exc^n(\mathcal{C},\mathcal{D})$. 
		Note that when $\mathcal{C}$ and $\mathcal{D}$ are both stable, the notion of excisive and coexcisive functors coincide. 
	\end{defn}

	As we are working with functors $F:\Sp_{T(n)}\to\Sp_{T(n)}$, we hence assume from now on $\mathcal{C}$ and $\mathcal{D}$ are both stable (even some arguments still work when they are not), so that the excisive functors are exactly the coexcisive functors. 

	\begin{const}\label{Cross_eff}
		We now introduce methods to detect whether a functor is $n$-excisive. We define the $n$-th cross effect of a functor $F:\mathcal{C}\to\mathcal{D}$ to be the functor $\mathcal{C}^n\to\mathcal{D}$ given by 
		\[ \crs_n(F)(X_1,\dots,X_n)=\tfib_{T\subseteq S}(F(\coprod_{s\notin T}X_s)) \]
		where $S=\{1,\dots,n\}$. 

		Dually we define the cocross effect to be 
		\[ \crs^n(F)(X_1,\dots,X_n)=\tcof_{T\subseteq S}(F(\prod_{s\in T}X_s)) \]

		The two functors detect whether $F$ sends strongly coCartesian $S$-cubes to Cartesian $S$-cubes (or its dual version), hence gives the obvious criterion that $F$ is $n$-excisive iff $\crs_{n+1}(F)=0$. 
		Note that the construction of (co)cross effects can be iterative, that is, we have 
		\[ \crs_n(F)(X_1,\dots,X_n)\simeq\crs_2(\crs_{n-1}(F)(X_1,\dots,X_{n-2},-))(X_{n-1},X_n) \]
		and 
		\[ \crs^n(F)(X_1,\dots,X_n)\simeq\crs^2(\crs^{n-1}(F)(X_1,\dots,X_{n-2},-))(X_{n-1},X_n) \]

		This follows from the iterative construction of the total (co)fibers. 
	\end{const}

	\begin{const}\label{Univ_exc_homog}
		We note that taking the cross effect of a functor $F:\mathcal{C}\to\mathcal{D}$ defines a functor
		\[ :\Fun(\mathcal{C},\mathcal{D})\to\Fun(\mathcal{C},\mathcal{D}) \]
		given by 
		\[ \crs_n^\Delta(F)(X)=\crs_n(F)(X,\dots,X) \]

		Since $\crs_n(X_1,\dots,X_n)$ can be written as a retract of $\crs_n(\coprod X_i,\dots,\coprod X_i)$, we have $F$ $n$-excisive iff $\crs_{n+1}^{\Delta}(F)\simeq*$. Noticing that $F\simeq\crs_1^\Delta(F)$, we have a natural transformation $\varepsilon:\crs_n^\Delta(F)\to F$ given by 
		\[ \crs_n\Delta(F)(X)\simeq\crs_n(F)(X,\dots,X)\to F(\coprod X)\to F(X) \]
		where the second map is given by the structure map of the total fiber. We define the cofiber of this natural tranformation to be $G_nF=\cof(\varepsilon)$, and define 
		\[ P_nF=\ilim(F\to G_{n+1}F\to G_{n+1}G_{n+1}F\to\cdots) \]

		By checking definitions we see that $P_nF$ is the couniversal $n$-excisive functor equipped with a natural transformation $F\to P_nF$. In other terms, we have $P_n$ a localization functor left adjoint to the inclusion $\Exc_n(\mathcal{C},\mathcal{D})\inj\Fun(\mathcal{C},\mathcal{D})$. 

		Dually we define 
		\[ \crs^n_\Delta(F)(X)=\crs^n(F)(X,\dots,X) \]
		and a functor is $n$-coexcisive iff $\crs^n_\Delta(F)\simeq*$. Also we have a natural transformation $\eta:F\to\crs^n_\Delta(F)$ given by 
		\[ F(X)\to F(\prod X)\to\crs^n(F)(X,\dots,X)\simeq\crs^n_\Delta(F)(X) \]
		take $G^nF=\fib(\eta)$, and finally 
		\[ P^nF=\plim(F\from G^{n+1}F\from\cdots) \]

		Again we have $P^nF$ is the universal $n$-coexcisive functor equipped with a natural transformation $P^nF\to F$. $P^n$ is also the colocalization functor right adjoint to the inclusion $\Exc^n(\mathcal{C},\mathcal{D})\to\Fun(\mathcal{C},\mathcal{D})$. 
		For a functor $F$, we define its $n$-homogeneous part $D_nF=\fib(P_nF\to P_{n-1}F)$. Dually its $n$-cohomogeneous part $D^nF=\cof(P^{n-1}F\to P^nF)$. 
	\end{const}

	\begin{defn}\label{Multilinear_fun}
		We define a functor $F:\mathcal{C}^n\to\mathcal{D}$ to be multilinear if it is 1-excisive in each variable. Dually we have the definition of comultilinear functors. 
	\end{defn}

	\begin{prop}\label{Exc_cross_eff_var}
		If $F:\mathcal{C}\to\mathcal{D}$ is $n$-excisive, then for $0\leq m\leq n$ we have the functor $\crs_{m+1}(F)$ is $(n-m)$-excisive in each variable. 
	\end{prop}
	\begin{proof}
		By the iterative construction of cross effects in \cref{Tot_fib}, we have 
		\[ \crs_{m+1}(F)(X_1,\dots,X_n,Y)\simeq\crs_m(F(-\coprod Y))(X_1,\dots,X_n) \] 
		and hence by induction it suffices to prove that if $F$ is $n$-excisive, then $F(-\coprod Y)$ is $(n-1)$-excisive, which is obvious. 
	\end{proof}

	\begin{thm}\label{N_homog_sym_multilinear}
		We have an equivalence of categories between the category of $n$-homogeneous functors $\mathcal{C}\to\mathcal{D}$ with the category of symmetric multilinear functors $\mathcal{C}^n\to\mathcal{D}$, with the equivalence given by $F(X)\mapsto\crs_n(F)(X_1,\dots,X_n)$ and $G(X_1,\dots,X_n)\mapsto G(X,\dots,X)_{hS_n}$. We define this second operation by $\Delta_{S_n}$. 

		Dually, we have an equivalence of categories of $n$-cohomogeneous functors with the category of symmetric multilinear functors, with the equivalence given by $F(X)\mapsto\crs^n(F)(X_1,\dots,X_n)$ and $G(X_1,\dots,X_n)\mapsto G(X,\dots,X)^{hS_n}$. We define this by $\Delta^{S_n}$. 
	\end{thm}
	\begin{remark}
		Note that in the stable settings, the $n$-excisive and coexcisive functors agree with each other, and the difference of $P_n$ and $P^n$ comes froms the fact that they are respectively the left and right adjoint to the inclusion. Here, the difference is similar. 
	\end{remark}
	\begin{proof}
		We first verify that for a symmetric multilinear functor $L$ we have $\crs_n\Delta_{S_n}L\to L$ is an isomorphism. By unwinding definitions, we have
		\begin{align*}
			&\crs_n\Delta_{S_n}L(X_1,\dots,X_n) \\
			\simeq&\tfib(T\mapsto L(\coprod_{s\notin T}X_s,\dots,\coprod_{s\notin T}X_s)_{hS_n}) \\
			\simeq&\tfib(T\mapsto L(\coprod_{s\notin T}X_s,\dots,\coprod_{s\notin T}X_s))_{hS_n} (\text{by~stability})
		\end{align*}

		Since $L$ is 1-excisive in each factor, we have 
		\[ L(\coprod_{s\notin T}X_s,\dots,\coprod_{s\notin T}X_s)\simeq\prod_{f\in T^n}L(X_{f(1)},\dots,X_{f(n)}) \]
		where $T\subseteq S$, $S=\{1,\dots,n\}$. If we denote this cube by $\mathcal{X}$, then this cube can be decomposed into the following. Let $\sigma:S\to S$ be an arbitrary map, and let $\mathcal{X}_\sigma$ denote the cube where $\mathcal{X}_\sigma(T)=L(X_{\sigma(1)},\dots,X_{\sigma(n)})$ if $\sigma(S)\subseteq S-T$ and is trivial otherwise. Then we clearly have $\mathcal{X}\simeq\prod_{\sigma:S\to S}\mathcal{X}_{\sigma}$. 

		Then note that if $\sigma$ is not a bijection we have $\tfib(\mathcal{X}_\sigma)$, then we have $\tfib(\mathcal{X}_\sigma)$ is trivial since in this cube, for $s\notin\img(\sigma)$, the morphism pointing from $T\to T\cup\{s\}$ is an isomorphism, and by the iterative construction of total fibers we have the total fiber is trivial by first taking fibers in this direction. 

		Finally for $\sigma$ is indeed a bijection, we have $\mathcal{X}_{\sigma}$ is trivial except for the initial vertex with value $L(X_{\sigma(1)},\dots,X_{\sigma(n)})$, hence we have 
		\[ \tfib(\mathcal{X}_\sigma)\simeq L(X_{\sigma(1)},\dots,X_{\sigma(n)}) \]

		Therefore 
		\[ \tfib(\mathcal{X})\simeq\prod_{\sigma\in S_n}L(X_{\sigma(1)},\dots,X_{\sigma(n)}) \]
		and the map $\crs_n\Delta_{S_n}L\to L$ is tautological. 

		Since $L$ is symmetric, we have 
		\[ \crs_n\Delta_{S_n}L(X_1,\dots,X_n)\simeq(\prod_{\sigma\in S_n}L(X_{\sigma(1)},\dots,X_{\sigma(n)}))_{hS_n}\simeq L(X_1,\dots,X_n) \]

		Then we treat the converse case. For an $n$-homogeneous functor $F$, note that $\crs_n(F)\simeq*$ implies that $F$ is also trivial, since $F$ is also $(n-1)$-excisive by definition. Therefore to prove that $\Delta_{S_n}\crs_n(F)\to F$ is an equivalence of $n$-homogeneous functors, it suffices to prove $\crs_n\Delta_{S_n}\crs_n(F)\to\crs_n(F)$ is an equivalence of symmetric multilinear functors. 

		We first construct the natural transformation $\gamma:\Delta_{S_n}\crs_n(F)\to F$. We define 
		\begin{align*}
			\gamma_X:&\Delta_{S_n}\crs_n(F)(X) \\
			\simeq&(\crs_n(F))(X,\dots,X)_{hS_n} \\
			\to&F(X\amalg\cdots\amalg X)_{hS_n} \\
			\to&F(X)_{hS_n} \\
			\to&F(X)
		\end{align*}
		where the first arrow is given by the structure map of the total fiber, and we have the last arrow since $F(X)$ has the trivial $S_n$-action. Then note that this map is an equivalence after passing to $\crs_n$ since it agrees with the equivalence defined above, we have the result. 
	\end{proof}
	
	\begin{coro}\label{Alt_const_n_homog}
		We have 
		\[ D_n(F)\simeq\ilim_m\Omega^{mn}(\crs_n(F))(\Sigma^mX,\dots,\Sigma^mX) \] 
		Dually we have 
		\[ D^n(F)\simeq\plim_m\Sigma^{mn}(\crs^n(F))(\Omega^mX,\dots,\Omega^mX) \]
	\end{coro}
	\begin{proof}
		Similar to the above argument. 
	\end{proof}

	\begin{coro}\label{Fil_colim_n_homog}
		A filtered-colimit preserving functor $F:\mathcal{C}\to\mathcal{D}$ where $\mathcal{C}$ is either compactly generated by $\mathbb{S}$ or stable and generated by $\Sigma^n\mathbb{S}$ for some object $\mathbb{S}$ is $n$-homogeneous iff it is of the form $F(X)=(C\otimes X^{\otimes n})_{hS_n}$. 
	\end{coro}
	\begin{proof}
		Using the above equivalence, we have $F(X)\simeq\crs_n(F)(X,\dots,X)_{hS_n}$. Therefore 
		\[ F(X)\simeq(\crs_n(F)(\mathbb{S},\dots,\mathbb{S})\otimes X^{\otimes n})_{hS_n} \]
		since $F$ hence $\crs_n(F)$ commutes with filtered colimits. 
	\end{proof}

	\begin{lemma}\label{n_homog_fac_stab}
		For any filtered colimit preserving $n$-homogeneous ($n\geq1$) functor $F:\mathcal{C}\to\mathcal{D}$, there is some $F':\Sp(\mathcal{C})\to\Sp(\mathcal{D})$ such that $F=\Omega^\infty F'\Sigma^\infty$. 
	\end{lemma}
	\begin{proof}
		By \cref{Alt_const_n_homog}, the image of a $n$-homogeneous functor lie in the subcategory of infinite loopspaces, and hence $F$ factors through $\Omega^\infty:\Sp(\mathcal{D})\to\mathcal{D}$. Now since $F$ factors as $\mathcal{C}\to\Sp(\mathcal{D})\to\mathcal{D}$, it also factors through $\Sigma^\infty:\mathcal{C}\to\Sp(\mathcal{C})$, and it suffices to prove that the middle functor remains $n$-homogeneous. This is direct by verifying the cross effects vanishes. 
	\end{proof}

	\begin{thm}[Telescopic Tate vanishing]\label{Tate_vanish}
		For $G$ a finite group, and an object $X$ of $\Sp_{T(n)}$ equipped with a $G$-action, then we have a norm map between the homotopy orbits and the homotopy fixed points $\Nm:X_{hG}\to X^{hG}$, and the cofiber of the norm map $X^{tG}=\cof(\Nm)$ vanishes. 
	\end{thm}
	\begin{proof}
		\todo{Reference}. 
	\end{proof}

	\begin{const}\label{N_exc_fun_char}
		Consider stable categories $\mathcal{C}$ and $\mathcal{D}$, a reduced $(n-1)$-excisive functor $E$, an $n$-homogeneous functor $K$, and an extension of $K$ into $E$ (that is, a fiber sequence $K\to F\to E$, which is also equivalent to the data $E\to\Sigma K$). 

		By \cref{N_homog_sym_multilinear}, we have $\Sigma K\simeq\Delta_{S_n}\crs_n(\Sigma K)$. Since $X$ is a retract of $X\oplus Y$, we have $F(X)$ a retract of $F(X\oplus Y)$, hence since $\mathcal{D}$ is also stable we have $F(X\oplus Y)\simeq F(X)\oplus K$. Thus $\cof(F(X)\to F(X\oplus Y))\simeq\fib(F(X\oplus Y)\to F(X))$. Thus by the direct definition of cross effects and the iterative algorithm for total (co)fibers we have the cross effect of $\Sigma K$ and the cocross effect agrees up to suspensions. Thus we have $\Delta^{S_n}(\crs^n(\Sigma K))$ is cohomogeneous, hence every natural transformation from an $(n-1)$-excisive functor $E$ to $\Delta^{S_n}(\Sigma K)$ vanishes. Finally by the long cofiber sequence 
		\[ (\crs_n(K)\Delta)^{tS_n}\to\Delta_{S_n}\crs_n(\Sigma K)\xrightarrow{\Nm}\Delta^{S_n}\crs_n(\Sigma K)\to(\crs_n(\Sigma K)\Delta)^{tS_n} \]
		and that the third functor in the cofiber sequence is $n$-cohomogeneous, we have that a map $E\to\Sigma K\simeq\Delta_{S_n}(\crs_n(\Sigma K))$ is equivalent to a map $E\to(\crs_n(K)\Delta)^{tS_n}$. 

		As a conclusion, a reduced $n$-excisive functor is equivalent to the data of a reduced $(n-1)$-excisive functor $E$, an $n$-homogeneous functor $K$ regarded as an $n$-symmetric functor, and a natural transformation $E\to(\crs_n(K)\Delta)^{tS_n}$. Moreover, this functor can be reconstructed as the fiber 
		\[ \fib(E\to(\crs_n(K)\Delta)^{tS_n}\to\Delta_{S_n}\crs_n(\Sigma K)\simeq\Sigma K) \]
	\end{const}

	\begin{coro}\label{Tel_loc_k_exc_split}
		Every $k$-excisive functor $F:\Sp_{T(n)}\to\Sp_{T(n)}$ splits as the direct sum of $j$-homogeneous functors (where $j<k$). 
	\end{coro}
	\begin{proof}
		Follows directly form \cref{Tate_vanish} and \cref{N_exc_fun_char}. 
	\end{proof}

\section{Dual calculus in the $T(n)$-local category and nilpotence}

	\cref{Tel_loc_k_exc_split} in the previous section implies that the Goodwillie tower of an $n$-excisive functor of $T(n)$-local spectra splits, which implies that the Goodwillie tower of a coanalytic functor splits. In this section, we wish to see that the converse result, that is, every split Goodwillie tower result in a coanalytic functor, then clearly it suffices to prove the following result. 

	\begin{thm}[Mathew]\label{Exc_comm_inf_prod}
		If $F\simeq\bigoplus F_j$, and $F_j:\Sp_{T(n)}\to\Sp_{T(n)}$ is filtered colimit preserving and $j$-homogeneous, then the natural map 
		\[ P^k(\bigoplus_{j=1}^\infty F_j)\simeq P^kF \]
		is an equivalence. 
		Dually we have 
		\[ P_k(\prod_{j=1}^\infty F_j)\simeq\prod_{j=1}^kF_j \]
		(though this dual is not needed later, we still present it here for the sake of symmetry). 
	\end{thm}

	We first present the core lemma of uniform nilpotence used to prove the theorem. 
	\begin{lemma}\label{Unif_nilp_homog}
		Let $F$ be a finite type $n$-spectrum. Then there exists a constant $C$ such that for all $j>k$ and any filtered colimit preserving $j$-homogeneous functor $H:\Sp_{T(n)}\to\Sp_{T(n)}$, we have the natural transformation 
		\[ F\otimes(\Sigma^{Ck}\crs^k_\Delta(H)\Omega^C)\to F\otimes(\crs^k_\Delta(H)) \]
		is null. This constant $C$ only depends on $F$. 
	\end{lemma}

	\begin{proof}
		We first show that \cref{Unif_nilp_homog} implies \cref{Exc_comm_inf_prod}. We first note that \cref{Exc_comm_inf_prod} is equivalent to $D^k(\bigoplus_{j>k} F_j)\simeq0$ for all $k$. Then we see that it suffices to prove that $D^k$ commutes with infinite direct sums, which is a property that ordinary sequential limits do not possess (for the construction of $D^k$ as a sequential limit, see \cref{Alt_const_n_homog}). But if we have \cref{Unif_nilp_homog}, this would give the fact that the inverse system $V\otimes\bigoplus(\Sigma^{Ck}\crs^k(F_j)\Delta\Omega^{C})$ is nilpotent. Furthermore, since $V$ is a finite spectrum, we have $V\otimes-$ commutes with sequential limits. Therefore 
		\[ 0\simeq\plim V\otimes\bigoplus(\Sigma^{nk}\crs^k(F_j)\Delta\Omega^{n})\simeq V\otimes\plim\bigoplus(\Sigma^{nk}\crs^k(F_j)\Delta\Omega^{n}) \]
		Then by taking the telescopic localization of $V$, we have $\plim\bigoplus(\Sigma^{nk}\crs^k(F_j)\Delta\Omega^{n})$ is $T(n)$-equivalent to $0$, which proves \cref{Exc_comm_inf_prod}. 
	\end{proof}

	Now it suffices to prove the lemma. To prove this, we need the language of nilpotence and apply it to naive equivariant homotopy theory. 

	\begin{defn}\label{Thick_tensor_ideal}
		We define a thick tensor ideal $\mathcal{T}^\otimes$ of $\mathcal{C}$ to be a thick subcategory which is also closed under tensoring with any object in $\mathcal{C}$, that is, for $X\in\mathcal{T}^\otimes$ and $Y\in\mathcal{C}$, $X\otimes Y$ is also in $\mathcal{T}^\otimes$. 
	\end{defn}

	\begin{const}\label{Thick_tensor_ideal_gen}
		For a symmetric monoidal category $\mathcal{C}$, recall the definition and the process of generation of thick subcategories in \cref{Thick_subcat} and \cref{Thick_subcat_gen}. 

		For a collection of objects $\mathcal{F}\subseteq\mathcal{C}$, since the intersection of two thick tensor ideals is still a thick tensor ideal, we may again consider the minimal thick tensor ideal containing $\mathcal{F}$, denoted by $\mathcal{T}^\otimes(\mathcal{F})$. We will call it the thick tensor ideal generated by $\mathcal{F}$, constructed as follows. 

		Let $\mathcal{T}^\otimes(\mathcal{F})=\mathcal{T}(\mathcal{F}^\otimes)$, where $\mathcal{F}^\otimes$ is the collections of objects in $\mathcal{C}$ of the form $X\otimes Y$, where $X\in\mathcal{C}$ and $\Sigma^nY\in\mathcal{F}$ for some $n\in\mathbb{Z}$. We can easily check that this construction satisfies the desired properties. Moreover, the filtration $\mathcal{T}_n(\mathcal{F})$ of $\mathcal{T}(\mathcal{F})$ induces a filtration on $\mathcal{T}^\otimes(\mathcal{F})$ by letting $\mathcal{T}^\otimes_n(\mathcal{F})=\mathcal{T}_n(\mathcal{F}^\otimes)$. By definition, each $\mathcal{T}^\otimes_n(\mathcal{F})$ is a tensor ideal of $\mathcal{C}$. 
	\end{const}

	\begin{defn}\label{Nilpotent_objects}
		For an object $A\in\mathcal{C}$, we define the subcategory of $A$-nilpotent objects to be $\mathcal{T}^\otimes(A)$, or equivalently, $\mathcal{T}(A\otimes\mathcal{C})$. We say it is $A$-nilpotent of degree $n$ if $n$ is the minimal natural number such that it lies in $\mathcal{T}_n^\otimes(A)$. Clearly, every nilpotent object is nilpotent of degree $n$ for some $n$. 
	\end{defn}

	\begin{lemma}\label{Nilp_equiv_defn}
		Let $I=\fib(1\to A)$. Then for an object $X\in\mathcal{C}$, the following are equivalent: 
		\begin{enumerate}
			\item $X\in\mathcal{T}_n^\otimes(A)$. 
			\item $X\otimes I^{\otimes n}\to X$ is null. 
			\item For every sequence of maps $X_n\to X_{n-1}\to\cdots\to X_0\simeq X$ with $A\otimes X_k\to A\otimes X_{k-1}$ null, the composition $X_n\to\cdots\to X$ is null. 
		\end{enumerate}
	\end{lemma}
	\begin{proof}
		$(1)\implies(2)$: For $n=0$, the assertion is trivial. If $n=1$, then we have the following retract diagram 
		\[ X\xrightarrow{f}A\otimes Y\xrightarrow{g}X \]
		Thus since $X\otimes I\to X\to X\otimes A$ is a fiber sequence, it suffices to prove $X\to X\otimes A$ admit a section, and hence $X\otimes A\simeq X\oplus\Sigma(X\otimes I)$, which implies $X\otimes I\to X$ is trivial. For this property, note 
		\[ X\xrightarrow{\eta\otimes\id}A\otimes X\xrightarrow{\id\otimes f}A\otimes A\otimes Y\xrightarrow{\mu\otimes 1} A\otimes Y \]
		is equivalent to $f$ by definition, since $\eta:1\to A$ is the unit map and $\mu:A\otimes A\to A$ is the multiplication map. Therefore, 
		\[ X\xrightarrow{\eta\otimes\id}A\otimes X\xrightarrow{\id\otimes f}A\otimes A\otimes Y\xrightarrow{\mu\otimes 1} A\otimes Y\xrightarrow{g}X \]
		exhibits $X$ as a retract of $A\otimes X$. For the general case, we apply induction and suppose that the case $n$ is solved. Then for the case of $n+1$, suppose $X$ is $A$-nilpotent of exponent $n+1$, let $Y\to X\to Z$ be a fiber sequence where $Y$ is $A$-nilpotent of exponent $1$ and $Z$ is $A$-nilpotent of exponent $n$. Then we have the following diagram where the columns are fiber sequences 
		% https://q.uiver.app/#q=WzAsOSxbMCwxLCJYXFxvdGltZXMgSV57XFxvdGltZXMgbisxfSJdLFswLDAsIllcXG90aW1lcyBJXntcXG90aW1lcyBuKzF9Il0sWzAsMiwiWlxcb3RpbWVzIElee1xcb3RpbWVzIG4rMX0iXSxbMiwxLCJYIl0sWzIsMiwiWiJdLFsyLDAsIlkiXSxbMSwwLCJZXFxvdGltZXMgSSJdLFsxLDEsIlhcXG90aW1lcyBJIl0sWzEsMiwiWlxcb3RpbWVzIEkiXSxbMSwwXSxbMCwyXSxbMyw0XSxbNSwzXSxbMSw2XSxbNiw1LCIwIl0sWzYsN10sWzcsOF0sWzIsOCwiMCIsMl0sWzgsNF0sWzAsN10sWzcsM11d
		\[\begin{tikzcd}[ampersand replacement=\&]
			{Y\otimes I^{\otimes n+1}} \& {Y\otimes I} \& Y \\
			{X\otimes I^{\otimes n+1}} \& {X\otimes I} \& X \\
			{Z\otimes I^{\otimes n+1}} \& {Z\otimes I} \& Z
			\arrow[from=1-1, to=1-2]
			\arrow[from=1-1, to=2-1]
			\arrow["0", from=1-2, to=1-3]
			\arrow[from=1-2, to=2-2]
			\arrow[from=1-3, to=2-3]
			\arrow[from=2-1, to=2-2]
			\arrow[from=2-1, to=3-1]
			\arrow[from=2-2, to=2-3]
			\arrow[from=2-2, to=3-2]
			\arrow[from=2-3, to=3-3]
			\arrow["0"', from=3-1, to=3-2]
			\arrow[from=3-2, to=3-3]
		\end{tikzcd}\]
		Then note that the map $X\otimes I^{\otimes n+1}\to Z\otimes I$ is null, hence lifts to a map $X\otimes I^{\otimes n+1}\to Y\otimes I$. Therefore, the map $X\otimes I^{\otimes n+1}\to X$ factors as 
		\[ X\otimes I^{\otimes n+1}\to Y\otimes I\xrightarrow{0}Y\to X \]
		hence is null. The case where $X'$ is the retract of the above $X$ is direct. 

		$(2)\implies(1)$: Suppose $\psi:X\otimes I^{\otimes n}\to X$ is null. Then $X$ is a retract of $\cof\psi$, hence it suffices to prove $\cof\psi\in\mathcal{T}_n^\otimes(A)$. To prove this, it suffices to prove $\cof(I^{\otimes n}\to 1)\in\mathcal{T}_n^\otimes(A)$, which is the extension of 
		\[ \cof(I^{\otimes k+1}\to I^{\otimes k})\simeq A\otimes I^{\otimes k}\in\mathcal{T}_1^\otimes(A) \]

		$(2)\implies(3)$: consider the following diagram 
		% https://q.uiver.app/#q=WzAsNSxbMCwxLCJYXzEiXSxbMSwxLCJYIl0sWzAsMiwiWF8xXFxvdGltZXMgQSJdLFsxLDIsIlhcXG90aW1lcyBBIl0sWzEsMCwiWFxcb3RpbWVzIEkiXSxbMCwxXSxbMSwzXSxbMCwyXSxbMiwzXSxbNCwxXSxbMCw0LCIiLDEseyJzdHlsZSI6eyJib2R5Ijp7Im5hbWUiOiJkYXNoZWQifX19XV0=
		\[\begin{tikzcd}[ampersand replacement=\&]
			\& {X\otimes I} \\
			{X_1} \& X \\
			{X_1\otimes A} \& {X\otimes A}
			\arrow[from=1-2, to=2-2]
			\arrow[dashed, from=2-1, to=1-2]
			\arrow[from=2-1, to=2-2]
			\arrow[from=2-1, to=3-1]
			\arrow[from=2-2, to=3-2]
			\arrow[from=3-1, to=3-2]
		\end{tikzcd}\]
		since the bottom map is null, the map $X_1\to X\to X\otimes A$ is null, hence lifts to $X_1\to X\otimes I$. Likewise, we may lift the composition $X_2\to X_1\to X\otimes I$ to $X_2\to X\otimes I^{\otimes2}$ by the same process. Inductively, we result in the fact that the map $X_n\to X$ factors as $X_n\to X\otimes I^{\otimes n}\to X$, hence is null by assumption. 

		$(3)\implies(2)$: Note that the following sequence 
		\[ X\otimes I^{\otimes n}\to X\otimes I^{\otimes n-1}\to\cdots\to X \]
		satisfy the desired assumption. 
	\end{proof}

	Then to apply the above definitions to equivariant settings, we apply another version of the monadicity theorem. 

	\begin{thm}\label{Beck_monadicity_tensor_cat}
		Let $F:\mathcal{C}\adj\mathcal{D}:G$ be an adjunction of idempotent complete symmetric monoidal stable categories, and $F$ is symmetric monoidal (this implies that $G$ is lax symmetric monoidal). Then if the counit $\varepsilon:FG\to\id_\mathcal{D}$ admit a section $\xi:\id_\mathcal{D}\to FG$, and that the projection formula $G(X)\otimes Y\simeq G(X\otimes F(Y))$ holds, then the adjunction is monadic, and the monad can be identified with the monad $G(1)\otimes-$ (here $1$ is the tensor unit of $\mathcal{D}$, and $G(1)$ admits a natural algebra structure since $G$ is lax symmetric monoidal). 
	\end{thm}
	\begin{proof}
		We first notice that the section of $\varepsilon$ gives a section of the monad $T=GF$, explicitly $G\xi F:T\to TT$ satisfies 
		\[ T\xrightarrow{G\xi F}TT\xrightarrow{\mu}T \]
		is the identity. 

		Then for any $T$-module $X$, note that the structure map $m:TX\to X$ admit a natural section $X\to TX$ given by the composition 
		\[ X\xrightarrow{\eta_X}TX\xrightarrow{(G\xi F)_X}T^2X\xrightarrow{Tm_X}TX \]
		This is a section of $m:TX\to X$ by definition of $G\xi F$. Moreover this natural section is $T$-linear, hence we have $X$ a retract of $TX$ for every $T$-module $X$. Therefore the inclusion $\LMod_T^{\mathrm{free}}(\mathcal{C})\inj\LMod_T(\mathcal{C})$ is an equivalence after applying idempotent completion, since it is fully faithful (by definition) and every object of $\LMod_T(\mathcal{C})$ is a retract of an object of $\LMod_T^{\mathrm{free}}(\mathcal{C})$. 

		We then prove that the map $K:\LMod_T^{\mathrm{free}}(\mathcal{C})\to\mathcal{D}$ (recall its construction in \cref{Adj_EM_K_fun_defn}) is an equivalence after applying idempotent completion. Firstly we have seen in \cref{Adj_EM_K_fun_defn} that $K$ is fully faithful. Then for an object of $\mathcal{D}$, the section $\xi$ provides 
		\[ D\xrightarrow{\xi_D}FGD\xrightarrow{\varepsilon_D}D \]
		hence we have $D$ a retract of $FGD\simeq KGD$. 

		Hence we have $E:\mathcal{D}\to\LMod_T(\mathcal{C})$ an equivalence after applying idempotent completion. But by assumption both sides are already idempotent complete, so the adjunction is monadic. 

		It then remains to identify the monads. Namely, by the projection formula, we have 
		\[ G(1)\otimes-\simeq G(1\otimes F-)\simeq GF- \]
		the identification of the units and multiplications follows by diagram chasing. 
	\end{proof}

	\begin{lemma}\label{Bialg_counit_section}
		Let $f:R\to S$ be a morphism of $\mathbb{E}_1$-rings. Suppose that there is a map of $R-R$-bimodules $g:S\to R$ such that $gf=1$, then the forgetful functor is left adjoint to $\Hom_R(S,-):\LMod_R\to\LMod_S$ ($\Hom_R(S,M)$ is equipped with the left $S$-module structure via right action on the first factor), and the counit admits a natural section. 
	\end{lemma}
	\begin{proof}
		Firstly the adjunction is obvious. Then the counit of the adjunction is given by 
		\[ f^*:\Hom_R(S,-)\to\Hom_R(R,-) \]
		Then we claim that $g^*$ is the section as desired. Firstly, note that 
		\[ g^*:\Hom_R(R,-)\to\Hom_R(S,-) \]
		is a morphism of left $R$-modules since $g:S\to R$ is a morphism of $R-R$-bimodules, and in particular preserves the right $R$-module structure. Then $f^*g^*=(gf)^*=\id$, so the counit admits a natural section. 
	\end{proof}

	\begin{thm}\label{Equivariant_subgrp_module_cat}
		For $G$ a finite group, let $H$ be a subgroup of $G$ (this is in fact true for stably dualizable group and a finite index subgroup), then we have $\Sigma^\infty_+G/H$ is a natural cocommutative coalgebra given by the diagonal map, and hence a natural commutative algebra structure on $\dual\Sigma^\infty_+G/H$. Then we have 
		\[ \Sp^{BH}\simeq\Mod_{\dual\Sigma^\infty_+G/H}(\Sp^{BG}) \]
		where $\Sp^{BG}$ is the category of naive $G$-spectra.  
	\end{thm}
	\begin{remark}
		This argument also applies to the category of genuine $G$-spectra with some care. 
	\end{remark}
	\begin{proof}
		We here regard $\Sp^{BG}$ as the category $\LMod_{\mathbb{S}[G]}(\Sp)$ (here $\mathbb{S}[G]:=\Sigma^\infty_+G$ is the $\mathbb{E}_1$-group algebra), equipped wit. Then note that the restriction functor is given by the forgetful functor of modules induced by the natural map $\mathbb{S}[H]\inj\mathbb{S}[G]$. The coinduction is then given by 
		\[ \coind_H^G(M)\simeq\Hom_{\mathbb{S}[H]}(\mathbb{S}[G],M) \]
		By \cref{Beck_monadicity_tensor_cat}, it suffices to show that the counit $\varepsilon:\res_H^G\circ\coind_H^G\to\id$ admits a section (the projection formula holds by formal reasons). But in fact, the section comes from \cref{Bialg_counit_section}, since the inclusion $\mathbb{S}[H]\inj\mathbb{S}[G]$ has obviously a $\mathbb{S}[H]-\mathbb{S}[H]$-bimodule which is a left inverse given by 
		\[ \mathbb{S}[G]\simeq\bigoplus_{g\in G}\mathbb{S}\to\bigoplus_{h\in H}\mathbb{S}\simeq\mathbb{S}[H] \]
		which is the identity on $\mathbb{S}$ indexed by $h\in H$ and null on $\mathbb{S}$ indexed by $g\notin H$. Therefore by \cref{Beck_monadicity_tensor_cat} this adjunction is monadic, and it is direct to verify the monad as the algebra $\Sigma^\infty_+G/H$. 
	\end{proof}

	\begin{const}\label{Subgrp_nilp_exp}
		We apply \cref{Equivariant_subgrp_module_cat} to nilpotence. Let $\mathcal{F}$ be a collection of subgroups of $G$. We say a object of $\Sp^{BG}$ is $\mathcal{F}$-nilpotent of degree $n$ if it is $\prod_{H\in\mathcal{F}}(\Sigma^\infty_+G/H)$-nilpotent of degree $n$. We will mainly apply the case where $\mathcal{F}$ consists of a single subgroup $H\leq G$. 
	\end{const}

	\begin{thm}\label{Fin_sp_nilp_fin_grp}
		Any finite spectrum $F$ in $\Sp_{T(n)}$ equipped with the trivial $G$-action is nilpotent viewed as an element of $\Sp^{BG}$. 
	\end{thm}
	\begin{proof}
		We consider the simplicial model for a point with free $G$-action, namely, let $EG_k=G^{k+1}$, and $EG=\ilim EG_{\leq k}$. Then since we are working in the $T(n)$-local category, by \cref{Tate_vanish} we have 
		\[ \mathbb{S}_{T(n)}^{hG}\simeq(\mathbb{S}_{T(n)})_{hG}\simeq(\ilim(\mathbb{S}_{T(n)}\otimes EG_{\leq k}))_{hG}\simeq\ilim(\mathbb{S}_{T(n)}\otimes EG_{\leq k})_{hG}\simeq\ilim(\mathbb{S}_{T(n)}\otimes EG_{\leq k})^{hG} \]
		By compactness of $F$, the map $F\inj F\otimes\mathbb{S}_{T(n)}^{hG}$ factors as 
		\[ F\to F\otimes(\mathbb{S}_{T(n)}\otimes EG_{\leq k})^{hG}\to F\otimes\mathbb{S}_{T(n)}^{hG} \]
		Since $F$ is finite, and homotopy fixed points agrees with the homotopy orbits, hence $F\otimes\mathbb{S}_{T(n)}^{hG}\simeq(F\otimes\mathbb{S}_{T(n)})^{hG}$, and the middle term satisfy a similar reduction. By the adjunction $\Delta\dashv\plim$ (here interpreted as the giving a spectrum the trivial $G$-action is left adjoint to taking homotopy fixed points), we obtain a retract diagram 
		\[ F\to F\otimes EG_{\leq k}\to F \]
		but $F\otimes EG_{\leq k}$ is clearly $G$-nilpotent by its construction, which implies that so is $F$. 
	\end{proof}

	\begin{lemma}\label{Unif_nilp_exp_fin_sp}
		For a finite spectrum $F$ (equipped with the trivial action), there is some constant $C$ only depending on $F$ such that for any finite group and a subgroup $H\leq G$ with $v_p(G)-v_p(H)\leq 1$ (here $v_p(G)$ is the $p$-adic valuation of the order of $G$), and any object $X\in\Sp_{T(n)}^{BG}$, we have $\exp_H(F\otimes X)\leq C$. 
	\end{lemma}
	\begin{proof}
		Consider the Sylow-$p$-subgroup $G_p\leq G$ and $H_p\leq H$, and we assume $H_p\leq G_p$ (which would imply $H_p\triangleleft G_p$ by some finite group theory). Consider $\coind_{G_p}^G(\dual\Sigma^\infty_+(G_p/H_p))=\dual\Sigma^\infty_+(G/H_p)$. Note that $G/H_p$ projects naturally to $G/H$, and by sending $gH$ to the sum of $ghH_p$ (where $hH_p$ runs over the cosets of $H_p$ in $H$), we obtain that $\dual\Sigma^\infty_+(G/H)$ is a retract of $\dual\Sigma^\infty_+(G/H_p)$. This implies that $\exp_{H_p}(\res_{G_p}^GX\otimes F)\geq\exp_H(X\otimes F)$, since if $\res_{G_p}^GX\otimes F\in\mathcal{T}^\otimes_n(\Sigma^\infty_+(G/H_p))$, by the same generation process, after the retraction, we would have $X\otimes F\in\mathcal{T}^\otimes_n(\Sigma^\infty_+(G/H))$. 

		Therefore we may pass to the case where $G$ is a $p$-group, and hence $H\leq G$ is an index $p$ subgroup of $G$, and hence normal. Therefore $\exp_H(F\otimes X)\leq\exp_H(F)$, and $\exp_H(F)\leq\exp(F)$, where the right hand side is the exponent of $F$ as a $C_p$-spectra. Therefore taking $C=\exp(F)$ gives the desired result. 
	\end{proof}

	Now we are ready to prove \cref{Unif_nilp_homog}
	\begin{proof}
		Firstly let $\mathbb{S}_{T(n)}^{\rho_j}$ denote the representation sphere of the standard representation of the symmetry group $S_j$ (that is, $\Sigma_j\act\mathbb{R}^{j-1}$). Then note that by the universal property of loopspaces, we have a $S_j$-equivariant map 
		\[ e:\mathbb{S}_{T(n)}\simeq\mathbb{S}_{T(n)}^{\otimes j}\to\Omega(\Sigma\mathbb{S}_{T(n)})^{\otimes j}\simeq\mathbb{S}_{T(n)}^{\rho_j} \]
		This is, in fact, the Euler class for this representation of $S_j$, which vanishes if the representation contains a trivial subrepresentation. We claim that $F\otimes e^C:F\to F\otimes\mathbb{S}_{T(n)}^{C\rho_j}$ is null-homotopic. This is since we may choose a Young subgroup $H\leq S_j$ with $v_p(S_j)-v_p(H)\leq 1$, and the restriction of the standard representation to any Young subgroup of $S_j$ admits a trivial subrepresentation. Therefore $e$ is null if restricted to $\Sp^{BH}$. By \cref{Nilp_equiv_defn}, we have $V\otimes e^C$ null in $\Sp^{BS_j}$ for $j\geq2$ (note $j\geq2$ is essential since otherwise the representation would be trivial anyway, and the Euler class would be the identity). 

		Then, we show that the natural transformation 
		\[ F\otimes(\Sigma^{Ck}\crs^k_\Delta(H)\Omega^C)\to F\otimes(\crs^k_\Delta(H)) \]
		is null. By the equivalence $H\simeq(L\circ\Delta)_{hS_j}$ as in \cref{N_homog_sym_multilinear} (where $L=\crs^j(H)$), it suffices to prove that the $S_j$-equivariant natural transformation 
		\[ F\otimes(\Sigma^{Ck}\crs^k_\Delta(L\circ\Delta_j\circ\Omega^C))\to F\otimes\crs^k_\Delta(L\circ\Delta_j) \]
		is null (the original case follows from taking $S_j$-homotopy orbits of this one). Since $L$ preserves filtered colimits in each variables by assumption, we have $L(X_1,\dots,X_j)\simeq\partial L\otimes X_1\otimes\cdots\otimes X_j$. By directly checking the definition of cocross effects, we have 
		\[ \Sigma^{Ck}\crs^k_\Delta(L\circ\Delta_j\circ\Omega^C)\simeq(\bigoplus_{f:[j]\surj[k]}\mathbb{S}_{T(n)}^{-C\rho_{f^{-1}(1)}}\otimes\cdots\otimes\mathbb{S}_{T(n)}^{-C\rho_{f^{-1}(k)}})\otimes(L\circ\Delta_j) \]
		as $S_j$-equivariant objects (where the $S_j$ action is given by the permutation on $[j]=\{1,2,\dots,j\}$, and in case that some $\sigma\in S_j$ preserves some surjection $f:[j]\surj[k]$, it must lie in the subgroup $\Sigma_{f^{-1}(1)}\times\cdots\times\Sigma_{f^{-1}(k)}$, and thus acts on the permutation spheres). Similarly, 
		\[ \crs^k_\Delta(L\circ\Delta_j)\simeq(\bigoplus_{f:[j]\surj[k]}\mathbb{S}_{T(n)}^{\otimes f^{-1}(1)}\otimes\cdots\otimes\mathbb{S}_{T(n)}^{\otimes f^{-1}(k)})\otimes(L\circ\Delta_j) \]
		where the action on $\mathbb{S}_{T(n)}$ is given likewise. Thus to show the desired equivariant map vanishes, it suffices to prove that for each Young subgroup $\Sigma_{j_1}\times\cdots\times\Sigma_{j_k}$, we have the equivariant map 
		\[ F\otimes\mathbb{S}_{T(n)}^{-C\rho_{j_1}}\otimes\cdots\otimes\mathbb{S}_{T(n)}^{-C\rho_{j_k}}\to F\otimes\mathbb{S}_{T(n)} \]
		vanishes. But taking dual of the above, we obtain 
		\[ \dual F\otimes e_1^C\otimes\cdots\otimes e_k^C:\dual F\to \dual F\otimes\mathbb{S}_{T(n)}^{C\rho_{j_1}}\otimes\cdots\otimes\mathbb{S}_{T(n)}^{C\rho_{j_k}} \]
		Now by the assumption $j>k$, we have at least one $j_i\geq2$, hence $\dual F\otimes\mathbb{S}_{T(n)}\to \dual F\otimes\mathbb{S}_{T(n)}^{C\rho_{j_i}}$ vanishes, hence the above map vanishes if we let $C$ be the constant as in \cref{Unif_nilp_exp_fin_sp} for $\dual F$.
	\end{proof}

	\begin{coro}\label{Coanalytic_fun_classification}
		Let $F:\Sp_{T(n)}\to\Sp_{T(n)}$ be a reduced functor preserving filtered colimits. Then $F$ is coanalytic (i.e. $\ilim P^kF\simeq F$) iff $F$ is of the form 
		\[ F(X)=\bigoplus(C_j\otimes X^{\otimes j})_{hS_j} \]
		where $C_j$ is an $S_j$-equivariant $T(n)$-local spectra, and the $S_j$-action on $X^{\otimes j}$ is given by permutation of the factors. 
	\end{coro}
	\begin{proof}
		$\implies$: 
		\[ F(X)\simeq\ilim_kP^kF(X)\simeq\ilim_k\bigoplus_{j<k}D_jP^kF(X)\simeq\ilim_k\bigoplus_j{j<k}(C(j)\otimes X^{\otimes j})\simeq\bigoplus_j(C(j)\otimes X^{\otimes j}) \]
		$\impliedby$: By \cref{Exc_comm_inf_prod}, 
		\[ D^j(\bigoplus)\simeq D^j(\bigoplus F_k)\simeq F_j \]
	\end{proof}
	
\section{Classical operads and deduction of the final statement}

	\begin{defn}\label{Classical_operad}
		We define a classical operad of a category $\mathcal{C}$ to be an $\mathbb{E}_1$-algebra in the category of symmetric sequences in $\mathcal{C}$, that is, $\Fun(\Core\Fin,\mathcal{C})$, with respect to the composition product. A classical cooperad is a $\mathbb{E}_1$-coalgebra with respect to the composition product. 
	\end{defn}

	\begin{const}\label{Tel_loc_monad_op}
		Note that the functor $\Fun(\Core\Fin,\mathcal{C})\to\Fun(\mathcal{C},\mathcal{C})$ given by 
		\[ F\mapsto (X\mapsto\bigoplus(F(n)\otimes X^{\otimes n})_{hS_n}) \]
		is symmetric monoidal, so any classical operad gives rise to a monad in $\mathcal{C}$. \cref{Coanalytic_fun_classification} shows that when $\mathcal{C}=\Sp_{T(n)}$, this inclusion is fully faithful with image the coanalytic functors, and $P^\infty=\ilim P^n$ is the associated coreflection functor (i.e. the right localization onto the coanalytic functors). 
	\end{const}

	\begin{const}\label{Bar_const}
		Recall that for any augmented $\mathbb{E}_1$-algebra of $\mathcal{C}$, say $f:A\to 1$, we define the bar construction of $A$ to be the realization of the simplicial object 
		% https://q.uiver.app/#q=WzAsNCxbMywwLCIxIl0sWzIsMCwiQSJdLFsxLDAsIkFcXG90aW1lcyBBIl0sWzAsMCwiXFxjZG90cyJdLFsxLDAsIiIsMCx7Im9mZnNldCI6LTF9XSxbMSwwLCIiLDIseyJvZmZzZXQiOjF9XSxbMiwxXSxbMiwxLCIiLDIseyJvZmZzZXQiOi0yfV0sWzIsMSwiIiwyLHsib2Zmc2V0IjoyfV0sWzMsMiwiXFxjZG90cyJdLFszLDIsIlxcY2RvdHMiLDJdXQ==
		\[\mBar(A)=\begin{tikzcd}[ampersand replacement=\&]
			\cdots \& {A\otimes A} \& A \& 1
			\arrow["\cdots", from=1-1, to=1-2]
			\arrow["\cdots"', from=1-1, to=1-2]
			\arrow[from=1-2, to=1-3]
			\arrow[shift left=2, from=1-2, to=1-3]
			\arrow[shift right=2, from=1-2, to=1-3]
			\arrow[shift left, from=1-3, to=1-4]
			\arrow[shift right, from=1-3, to=1-4]
		\end{tikzcd}\]
		By definition this is the relative tensor $1\otimes_A1$, and it admits a natural augmented $\mathbb{E}_1$-coalgebra structure given by 
		\[ 1\otimes_A1\simeq 1\otimes_AA\otimes_A1\to 1\otimes_A1\otimes_A1\simeq 1\otimes_A1\otimes1\otimes_A1 \]
		Similarly, for an augmented $\mathbb{E}_1$-coalgebra $B$ of $\mathcal{C}$, we may likewise define its cobar construction 
		% https://q.uiver.app/#q=WzAsNCxbMCwwLCIxIl0sWzEsMCwiQiJdLFsyLDAsIkJcXG90aW1lcyBCIl0sWzMsMCwiXFxjZG90cyJdLFswLDEsIiIsMix7Im9mZnNldCI6MX1dLFswLDEsIiIsMCx7Im9mZnNldCI6LTF9XSxbMSwyXSxbMSwyLCIiLDAseyJvZmZzZXQiOjJ9XSxbMSwyLCIiLDAseyJvZmZzZXQiOi0yfV0sWzIsMywiXFxjZG90cyIsMl0sWzIsMywiXFxjZG90cyJdXQ==
		\[\coBar(B)=\begin{tikzcd}[ampersand replacement=\&]
			1 \& B \& {B\otimes B} \& \cdots
			\arrow[shift right, from=1-1, to=1-2]
			\arrow[shift left, from=1-1, to=1-2]
			\arrow[from=1-2, to=1-3]
			\arrow[shift right=2, from=1-2, to=1-3]
			\arrow[shift left=2, from=1-2, to=1-3]
			\arrow["\cdots"', from=1-3, to=1-4]
			\arrow["\cdots", from=1-3, to=1-4]
		\end{tikzcd}\]
		This is likewise an augmented $\mathbb{E}_1$-algebra $\coBar(B)$ of $\mathcal{C}$. By definition, we have $\mBar:\Alg(\mathcal{C})\adj\co\Alg(\mathcal{C}):\coBar$. 
	\end{const}

	\begin{thm}\label{Susp_loop_adj_mon_coan}
		The functor $L_{T(n)}\Sigma^\infty_{T(n)}\Omega^\infty_{T(n)}:\Sp_{T(n)}\to\Sp_{T(n)}$ (recall this construction in \cref{Tel_susp_loop}) is coanalytic. More precisely, we have 
		\[ L_{T(n)}\Sigma^\infty_{T(n)}\Omega^\infty_{T(n)}X\simeq L_{T(n)}\bigoplus_{k=1}^\infty(X^\otimes k)_{hS_k} \]
	\end{thm}
	\begin{proof}
		For a $T(n)$-local spectrum $Y$, we apply $\Phi$ (the Bousfield-Kuhn functor) to the unit map $\Omega^\infty Y\to\Omega^\infty\Sigma^\infty\Omega^\infty Y$, and since $\Phi\Omega\simeq L_{T(n)}$ by \cref{Bousfield_Kuhn_tel_loc_fac}, we have a natural transformation $\lambda_Y:Y\to L_{T(n)}\Sigma^\infty\Omega^\infty Y$. Note that $L_{T(n)}\Sigma^\infty\Omega^\infty Y$ is a non-unital $\mathbb{E}_\infty$-ring spectrum, hence $\lambda$ extends to a map 
		\[ L_{T(n)}\bigoplus_{k=1}^\infty(Y^{\otimes k})\to L_{T(n)}\Sigma^\infty\Omega^\infty Y \]
		Now, by a result of Kuhn this is an equivalence when $X$ is $d_{n+1}$-connected (where $d_{n+1}$ is some constant depending on the inclusion of $\An_*^{v_n}\inj\An_*$, or the connectivity of some chosen finite space of type $n+1$) and $T(i)_*X=0$ for $0<i<n$. By \cref{Tel_susp_loop} and picking the corresponding inclusion of $\An_*^{v_n}$ into $\An_*$ we have alternatively $\Omega^\infty_{T(n)}$ equivalent to $\tau_{\geq d_{n+1}+1}\Omega^\infty:L_{T(n)}\Sp\to\An_*^{v_n}$, hence we have 
		\[ \Sigma^\infty_{T(n)}\Omega^\infty_{T(n)}X\simeq L_{T(n)}\Sigma^\infty\Omega^\infty(\tau_{\geq d+1}X) \]
		since $\tau_{\geq d+1}X$ is $T(n)$-local, it is $T(i)$-acyclic for $i<n$, and hence we have 
		\[ L_{T(n)}\Sigma^\infty_{T(n)}\Omega^\infty_{T(n)}(X)\simeq L_{T(n)}\bigoplus_{k=1}^\infty(\tau_{\geq d_{n+1}+1}(X)^{\otimes k})_{hS_k} \]
		Via the equivalence $L_{T(n)}\tau_{\geq d_{n+1}}X\simeq L_{T(n)}X$, the conclusion holds. 
	\end{proof}

	\begin{defn}\label{Cocomm_coop_Lie_operad}
		We define the cocommutative cooperad to be $\Sigma^\infty_{T(n)}\Omega^\infty_{T(n)}$ (this is coanalytic, and is also a comonad since it is the composition of a right adjoint by its left adjoint). We define the Lie operad to be its cobar construction. 
	\end{defn}

	\begin{remark}
		Algebras of Lie operads behaves like Lie algebras, and hence the final theorem would give a categorical characterization of Whitehead products. 
	\end{remark}

	\begin{thm}
		Let $F:\Sp_{T(n)}\to\Sp_{T(n)}$ be a functor that preserves sifted colimits. Then $F$ is coanalytic. 
	\end{thm}
	\begin{proof}
		We prove $F$ is a colimit of coanalytic functors, which is coanalytic since the inclusion of coanalytic functors into all functors preserves colimits. Note that for any spectrum $X$, the natural map 
		\[ \ilim_k\Sigma^{\infty-k}\Omega^{\infty-k}(\tau_{\geq d_{n+1}+1}X)\to X \]
		is an equivalence (here $d_{n+1}$ is defined as in the above proof, and is related to the explicit inclusion $\An_*^{v_n}\inj\An_*$. ), since this obviously holds for the case of a finite spectrum, and the general case follows from writing $X$ a filtered colimit of finite spectra. Therefore, it suffices to prove that the functor 
		\[ X\mapsto F(L_{T(n)}\Sigma^{\infty-k}\Omega^{\infty-k}(\tau_{\geq d_{n+1}+1}X)) \]
		is coanalytic. 

		Since $F$ preserves sifted colimits, and $L_{T(n)}$ and $\Sigma^{\infty-k}$ are left adjoints, we have $F\circ L_{T(n)}\circ\Sigma^{\infty-k}:\An_*\to\Sp_{T(n)}$ preserves sifted colimits. Since $\An_*$ is projectively generated with compact generators given by $\Fin_*$, $F\circ L_{T(n)}\circ\Sigma^{\infty-k}$ is the left Kan extension along the inclusion $i:\Fin_*\inj\An_*$ of its restriction $\Fin_*\to\Sp_{T(n)}$, which we denote by $f$. Now, recall that to calculate a left Kan extension, we may apply the coend formula, which is given as follows 
		\[ \Lan_if(X)\simeq\int^{I\in\Fin_*}\Hom_*(I,X)\otimes f(I) \]
		here this tensor is the tensor of a space with a spectrum. Recall also that we may reduce a coend to a colimit by passing to the twisted arrow category $T(\Fin_*)$ (that is, the straightening of $\Hom:\mathcal{C}\times\mathcal{C}^{\op}\to\An$), explicitly given by 
		\[ T(\Fin_*)\to\Fin_*^{\op}\times\Fin_*\xrightarrow{\Hom(-,X)\times f}\An_*\times\Sp_{T(n)}\xrightarrow{\otimes}\Sp_{T(n)} \]

		Thus, $F\circ L_{T(n)}\circ\Sigma^{\infty-k}$ can be written as the colimit of functors of the form 
		\[ \Hom_*(I,-)\otimes F(L_{T(n)}\Sigma^{\infty-k}J):\An_*\to\Sp_{T(n)} \]
		where $I$ and $J$ are finite pointed sets. This implies that 
		\[ X\mapsto F(L_{T(n)}\Sigma^{\infty-k}\Omega^{\infty-k}(\tau_{\geq d_{n+1}+1}X)) \]
		can be written as the colimit of functors of the form 
		\[ X\mapsto\Hom_*(I,\Omega^{\infty-k}(\tau_{\geq d_{n+1}+1}X))\otimes C \]
		for $C$ some $T(n)$-local spectrum. It suffices to show that 
		\[ X\mapsto L_{T(n)}\Sigma^\infty\Hom_*(I,\Omega^{\infty-k}(\tau_{\geq d_{n+1}+1}X)) \]
		are coanalytic. However, since $I$ is a finite set, the right hand side can be identified with $(\Sigma^\infty\Omega^{\infty-k}(\tau_{\geq d_{n+1}+1}X))^{\otimes I-*}$. But this functor is clearly coanalytic by the same process of proof as in \cref{Susp_loop_adj_mon_coan}. 
	\end{proof}

	\begin{thm}\label{Bousfield_Kuhn_Tel_loc_Lie_alg}
		There is a map of operads $\gamma:\Phi\Theta\to\coBar(\Sigma^\infty_{T(n)}\Omega^\infty_{T(n)})$ which is an equivalence of operads. Therefore, $\Phi\Theta$ can be identified with the Lie operad on $\Sp_{T(n)}$, which produces $\An_*^{v_n}$ as the category of Lie algebras of $\Sp_{T(n)}$. 
	\end{thm}
	\begin{proof}
		Note that the cobar construction $\coBar(\Sigma^\infty_{T(n)}\Omega^\infty_{T(n)})$ is given by the totalization of the cosimplicial object 
		% https://q.uiver.app/#q=WzAsNCxbMCwwLCJcXGlkX3tcXFNwX3tUKG4pfX0iXSxbMSwwLCJcXFNpZ21hXlxcaW5mdHlfe1Qobil9XFxPbWVnYV5cXGluZnR5X3tUKG4pfSJdLFsyLDAsIlxcU2lnbWFeXFxpbmZ0eV97VChuKX1cXE9tZWdhXlxcaW5mdHlfe1Qobil9XFxTaWdtYV5cXGluZnR5X3tUKG4pfVxcT21lZ2FeXFxpbmZ0eV97VChuKX0iXSxbMywwLCJcXGNkb3RzIl0sWzAsMSwiIiwyLHsib2Zmc2V0IjoxfV0sWzAsMSwiIiwwLHsib2Zmc2V0IjotMX1dLFsxLDJdLFsxLDIsIiIsMCx7Im9mZnNldCI6Mn1dLFsxLDIsIiIsMCx7Im9mZnNldCI6LTJ9XSxbMiwzLCJcXGNkb3RzIiwyXSxbMiwzLCJcXGNkb3RzIl1d
		\[\begin{tikzcd}[ampersand replacement=\&]
			{\id_{\Sp_{T(n)}}} \& {\Sigma^\infty_{T(n)}\Omega^\infty_{T(n)}} \& {\Sigma^\infty_{T(n)}\Omega^\infty_{T(n)}\Sigma^\infty_{T(n)}\Omega^\infty_{T(n)}} \& \cdots
			\arrow[shift right, from=1-1, to=1-2]
			\arrow[shift left, from=1-1, to=1-2]
			\arrow[from=1-2, to=1-3]
			\arrow[shift right=2, from=1-2, to=1-3]
			\arrow[shift left=2, from=1-2, to=1-3]
			\arrow["\cdots"', from=1-3, to=1-4]
			\arrow["\cdots", from=1-3, to=1-4]
		\end{tikzcd}\]
		But since $\Phi\Omega^\infty_{T(n)}\simeq\Sigma^\infty_{T(n)}\Theta\simeq\id_{\Sp_{T(n)}}$, the totalization of the above cosimplicial diagram is equivalent to the following % https://q.uiver.app/#q=WzAsMyxbMCwwLCJcXFBoaVxcT21lZ2FeXFxpbmZ0eV97VChuKX1cXFNpZ21hXlxcaW5mdHlfe1Qobil9XFxUaGV0YSJdLFsxLDAsIlxcUGhpXFxPbWVnYV5cXGluZnR5X3tUKG4pfVxcU2lnbWFeXFxpbmZ0eV97VChuKX1cXE9tZWdhXlxcaW5mdHlfe1Qobil9XFxTaWdtYV5cXGluZnR5X3tUKG4pfVxcVGhldGEiXSxbMiwwLCJcXGNkb3RzIl0sWzAsMSwiIiwyLHsib2Zmc2V0IjoxfV0sWzAsMSwiIiwwLHsib2Zmc2V0IjotMX1dLFsxLDJdLFsxLDIsIiIsMCx7Im9mZnNldCI6Mn1dLFsxLDIsIiIsMCx7Im9mZnNldCI6LTJ9XV0=
		\[\begin{tikzcd}[ampersand replacement=\&]
			{\Phi\Omega^\infty_{T(n)}\Sigma^\infty_{T(n)}\Theta} \& {\Phi\Omega^\infty_{T(n)}\Sigma^\infty_{T(n)}\Omega^\infty_{T(n)}\Sigma^\infty_{T(n)}\Theta} \& \cdots
			\arrow[shift right, from=1-1, to=1-2]
			\arrow[shift left, from=1-1, to=1-2]
			\arrow[from=1-2, to=1-3]
			\arrow[shift right=2, from=1-2, to=1-3]
			\arrow[shift left=2, from=1-2, to=1-3]
		\end{tikzcd}\]
		hence $\coBar(\Sigma^\infty_{T(n)}\Omega^\infty_{T(n)})\simeq\Tot(\Phi(\Omega^\infty_{T(n)}\Sigma^\infty_{T(n)})^{\bullet+1}\Theta)$, which produces a map 
		\[ \gamma:\Phi\Theta\to\Tot(\Phi(\Omega^\infty_{T(n)}\Sigma^\infty_{T(n)})^{\bullet+1}\Theta) \]
		Now upon applying $P_k$ on both sides, since $\Phi$ and $\Theta$ preserves filtered colimits, by the construction of $P_k$, we have $P_k(\Phi\Theta)\simeq\Phi P_k(\id_{\An_*^{v_n}})\Theta$, and on the right hand side we have 
		\[ P_k\Tot(\Phi(\Omega^\infty_{T(n)}\Sigma^\infty_{T(n)})^{\bullet+1}\Theta)\simeq\Phi\Tot(P_k((\Omega^\infty_{T(n)}\Sigma^\infty_{T(n)})^{\bullet+1}))\Theta \]
		Thus it suffices to verify that we have 
		\[ P_k(\id_{\An_*^{v_n}})\to\Tot(P_k((\Omega^\infty_{T(n)}\Sigma^\infty_{T(n)})^{\bullet+1})) \]
		is an equivalence. 

		In fact, we prove a more general statement saying that 
		\[ P_k(FG)\simeq\Tot(P_k(F(\Omega^\infty_{T(n)}\Sigma^\infty_{T(n)})^{\bullet+1}G)) \]
		Suppose $F$ is of the form $H\Sigma^\infty_{T(n)}$, then the above is obviously an equivalence by the extra degeneracy argument on cosimplicial objects. Then note that homogeneous functors are of the form $H\Sigma^\infty_{T(n)}$ by \cref{n_homog_fac_stab}. Therefore the assumption holds for all homogeneous functor $F$. Then the conclusion follows from induction and approximating $F$ to the $k$-th degree (note $P_k(P_k(F))=P_k(F)$), and on each step $D_iF\to P_iF\to P_{i-1}F$ is a fiber sequence, so the homogeneous case and the induction hypothesis works. 
	\end{proof}





%	\begin{proof}
%		This map originates from the property stated in \cref{Bousfield_Kuhn_susp_loop_adj}. We first note that in a general category $\mathcal{C}$, then a map of $\mathbb{E}_1$-algebras $A\to\coBar(B)$ is equivalent to the data of an augmented $\mathbb{E}_1$-algebra structure of some element projecting to $(A,B)$ in the twisted arrow category $T(\mathcal{C})$ (recall that this is the straightening of the functor $\Hom:\mathcal{C}\times\mathcal{C}^{\op}\to\An$). Therefore, to construct $\gamma$, it suffices to construct an augmented $\mathbb{E}_1$-algebra structure on the object $\Gamma:\Phi\Theta\to\Sigma^\infty_{T(n)}\Omega^\infty_{T(n)}$, where $\Gamma$ is given by 
%		\[ \Gamma:\Phi\Theta\xrightarrow{\Phi\eta^2\Theta}\Phi\Omega^\infty_{T(n)}\Sigma^\infty_{T(n)}\Omega^\infty_{T(n)}\Sigma^\infty_{T(n)}\Theta\simeq\Sigma^\infty_{T(n)}\Omega^\infty_{T(n)} \]
		
%		Then we construct an algebra structure on $\Gamma$. Note that for a symmetric monoidal action $\mathcal{C}\act\mathcal{D}$, an endomorphism object of $D\in\mathcal{D}$ is some $C\in\mathcal{C}$ equipped with a morphism $C\otimes D\to D$, such that it is universal in the sense that for all $E\in\mathcal{C}$, we have
%		\[ \Hom(E,C)\to\Hom(E\otimes D,C\otimes D)\to\Hom(E\otimes D,D) \]
%		is an equivalence. By Lurie, we have an $\mathbb{E}_1$-algebra structure on $C$. Since $\Fun(\Sp_{T(n)},\Sp_{T(n)})$ acts on $\Fun(\An_*^{v_n},\Sp_{T(n)})$ by composition of functors, we have $T(\Fun(\Sp_{T(n)},\Sp_{T(n)}))$ acts on $T(\Fun(\An_*^{v_n},\Sp_{T(n)}))$. Therefore, it suffices to show that $\Gamma$ is the endomorphism object of $\Phi\to\Sigma^\infty_{T(n)}$ (this is given by the same process as above). 

%		The action diagram is clearly given by 
		% https://q.uiver.app/#q=WzAsNCxbMCwwLCJcXFBoaVxcVGhldGFcXFBoaSJdLFsxLDAsIlxcU2lnbWFeXFxpbmZ0eV97VChuKX1cXE9tZWdhXlxcaW5mdHlfe1Qobil9XFxTaWdtYV5cXGluZnR5X3tUKG4pfSJdLFswLDEsIlxcUGhpIl0sWzEsMSwiXFxTaWdtYV5cXGluZnR5X3tUKG4pfSJdLFswLDIsIlxcdmFyZXBzaWxvbiIsMl0sWzMsMSwiXFxldGEiLDJdLFswLDFdLFsyLDNdXQ==
%		\[\begin{tikzcd}[ampersand replacement=\&]
%			{\Phi\Theta\Phi} \& {\Sigma^\infty_{T(n)}\Omega^\infty_{T(n)}\Sigma^\infty_{T(n)}} \\
%			\Phi \& {\Sigma^\infty_{T(n)}}
%			\arrow[from=1-1, to=1-2]
%			\arrow["\varepsilon"', from=1-1, to=2-1]
%			\arrow[from=2-1, to=2-2]
%			\arrow["\eta"', from=2-2, to=1-2]
%		\end{tikzcd}\]
%		Let $E:F\to G$ be an arbitrary object of $T(\Fun(\Sp_{T(n)},\Sp_{T(n)}))$, and we must show that 
%		\[ \Hom(F\to G,\Phi\Theta\to\Sigma^\infty_{T(n)}\Omega^\infty_{T(n)})\to\Hom(F\Phi\to G\Sigma^\infty_{T(n)},\Phi\to\Sigma^\infty_{T(n)}) \]
%		is an equivalence. The opposite of this map is given by 
%		\begin{align*}
%			&\Hom(F\Phi\to G\Sigma^\infty_{T(n)},\Phi\to\Sigma^\infty_{T(n)}) \\
%			\to&\Hom(F\Phi\Theta\to G\Sigma^\infty_{T(n)}\Omega^\infty_{T(n)},\Phi\Theta\to\Sigma^\infty_{T(n)}\Omega^\infty_{T(n)}) \\
%			\to&\Hom(F\to G,\Phi\Theta\to\Sigma^\infty_{T(n)}\Omega^\infty_{T(n)}) 
%		\end{align*}
%		and they are inverses to one another by the property shown in \cref{Bousfield_Kuhn_susp_loop_adj}. 

%		Finally it suffices to show that the map $\gamma$ induced by $\Gamma$ is an equivalence. Since the domain and codomain of $\gamma$ is coanalytic, it suffices to show that $\gamma$ induces an equivalence on the $k$-excisive approximations for all $k$. 
%	\end{proof}




%	\begin{const}\label{Nilpotence_exponent}
%		Let $G$ be a finite group, and we consider the category $\Fun(BG,\Sp_{T(n)})$. We explicitly construct the thick subcategory of $\Fun(BG,\Sp_{T(n)})$ spanned by free objects, that is, objects of the form $\Sigma^\infty G_+\otimes X$ with the natural $G$-action. We denote the collection of these objects by $\mathcal{F}$. Now, we define $\mathcal{T}_1(\mathcal{F})$ to be the collection of objects retracts of objects in $\mathcal{F}$. Then inductively, we define $\mathcal{T}_n(\mathcal{F})$ to be the collection of objects $X$ such that $X$ is a retract of some $X_0$, and $X_0$ sits in a fiber sequence $Y\to X_0\to Z$ where $Y\in\mathcal{T}_1(\mathcal{F})$ and $Z\in\mathcal{T}_{n-1}(\mathcal{F})$. 
%
%		Now we have $\bigcup\mathcal{T}_n(\mathcal{F})=\mathcal{T}(\mathcal{F})$ the thick subcategory generated by $\mathcal{F}$, and provides a filtration for this category. We say $X$ is nilpotent if $X\in\mathcal{T}(\mathcal{F})$, and define the exponent of $X$ $\exp(X)$ to be the minimal positive number $n$ such that $X\in\mathcal{T}_n(\mathcal{F})$. 
%		
%		Similarly for any subgroup $H\leq G$, we may define the notion of $H$-nilpotence and $H$-exponent for $H$-nilpotent objects $X$, denoted by $\exp_H(X)$. 
%	\end{const}


%	\begin{const}
%		Note that by definition, the functors $\crs_n^\Delta$ and $\crs^n_\Delta$ admit $S_n$-actions on them, namely by permuting the factors of $X$ in the slots of $\crs_n$ and $\crs^n$. This means that the natural transformations $\varepsilon:\crs_n^\Delta(F)\to F$ and $\eta:F\to\crs^n_\Delta(F)$ factors through the homotopy orbits and homotopy fixed points, namely we let 
%		\[ \varepsilon_n:\crs_n^\Delta(F)_{hS_n}\to F \]
%		and 
%		\[ \eta^n:F\to\crs^n_\Delta(F)^{hS_n} \]
%		be the maps factoring $\varepsilon$ and $\eta$. 
%	\end{const}

%	\begin{thm}[Goodwillie]\label{N_homog_fun_char}
%		For sufficiently good categories $\mathcal{C}$ and $\mathcal{D}$, we define $\mathcal{C}^{(n)}$ to be the category $\mathcal{C}^n_{hS_n}$, where the $S_n$ action on $\mathcal{C}^n$ is given by the permutation on the summands. We then define a symmetric functor $\mathcal{C}\to\mathcal{D}$ to be a functor $F:\mathcal{C}^{(n)}\to\mathcal{D}$. \\
%		Since the coproduct on $\mathcal{C}$ is symmetric monoidal, we have the coproduct map $\coprod:\mathcal{C}^n\to\mathcal{C}$ factors through the homotopy orbit of the $S_n$-action, hence produces a map $\coprod:\mathcal{C}^{(n)}\to\mathcal{C}$. Finally we define for a functor $F:\mathcal{C}\to\mathcal{D}$ its cross effect $\mathrm{cr}_{(n)}(F)$ to be its precomposition with $\coprod:\mathcal{C}^{(n)}\to\mathcal{C}$. \\
%		We have the following fully faithful embedding of categories 
%		\[ \mathrm{cr}_{(n)}:\mathrm{Homog}^n(\mathcal{C},\mathcal{D})\to\mathrm{SymFun}^n(\mathcal{C},\mathcal{D}) \]
%		Moreover it has essential image the functors whose underlying functor $F:\mathcal{C}^n\to\mathcal{D}$ that is multilinear (exact in each component). 
%	\end{thm}
%	\begin{proof}
%		\todo{Reference}. 
%	\end{proof}




%	(regarded as a set of objects of $\mathcal{C}$ indexed by the set of natural numbers, and equipped with an $S_n$-action for the object indexed by $n$)
%	\begin{defn}\label{Coanalytic_fun_defn_1}
%		We define a symmetric sequence $C$ in $\Sp_{T(n)}$ to be a collection of $T(n)$-local spectra $\{C(k)\}_{k\geq0}$ with $C(k)$ equipped with a $S_k$-action. A symmetric sequence $C$ gives rise to a functor $\Sp_{T(n)}\to\Sp_{T(n)}$ given by 
%		\[ \lambda(C):\Sp_{T(n)}\to\Sp_{T(n)},F_C(X)=\bigoplus_{k\geq0}(C(k)\otimes X^{\otimes k})_{hS_k} \]
%		Note that if we let $\Fin$ denote the category of finite sets, and $[n]$ be the set with $n+1$-elements, then we have $\Aut([n])\simeq BS_{n+1}$, where $S_{n+1}$ is the permuatation group of $n+1$-th order. Thus equivalently we may define the category of symmetric sequences to be $\Fun(\Core\Fin,\Sp_{T(n)})$. 
%	\end{defn}

%	The main purpose of this section is to prove the following theorem. 

%	\begin{thm}\label{Sym_seq_ff_emb}
%		The embedding $C\mapsto \lambda_C:\Fun(\Core\Fin,\Sp_{T(n)})\to\Fun(\Sp_{T(n)},\Sp_{T(n)})$ is fully faithful. We define its image to be the coanalytic functors. 
%	\end{thm}

%	To prove the above theorem, we need the following definitions. 

%	\begin{const}\label{Coder_fun_tel_loc}
%		We define a functor $\lambda_k:\Sp_{T(n)}^{BS_k}\to\Fun(\Sp_{T(n)},\Sp_{T(n)})$ given by 
%		\[ C\mapsto(X\mapsto (C\otimes X^{\otimes k})_{hS_k}) \]
%		where the $S_k$-action on $C\otimes X^{\otimes k}$ is given by the natural action on $C$ and the permutation action on $X^{\otimes k}$. Note that by construction $\lambda_k$ preserves all colimits, hence by the adjoint functor theorem it admits a right adjoint 
%		\[ \partial^k:\Fun(\Sp_{T(n)},\Sp_{T(n)})\to\Sp_{T(n)}^{BS_k} \]
%		Obviously to prove \cref{Sym_seq_ff_emb}, it suffices to prove that for any $C\in\Fun(\Core\Fin,\Sp_{T(n)})$, we have $C(k)\simeq\partial^k\lambda(C)$. 
%	\end{const}

%	\begin{lemma}\label{Char_coder_fun}
%		We have
%		\[ \partial^k(F)\simeq\Map_{\Fun(\Sp_{T(n)},\Sp_{T(n)})}((-)^{\otimes k},F) \]
%		Here note that $\Fun(\Sp_{T(n)},\Sp_{T(n)})$ is a stable category, so it makes sense to talk of $\Map(-,-)$ in this category (denoting the mapping spectrum). Moreover, the $S_k$-action of $\partial^k(F)$ is given by the inverse permutation action on $(-)^{\otimes k}$. 
%	\end{lemma}
%	\begin{proof}
%		We have the following chain of equivalences. 
%		\begin{align*}
%			&\partial^k(F) \\
%			\simeq&\Map_{\Sp_{T(n)}}(\mathbb{S}_{T(n)},\partial^k(F)) \\
%			\simeq&\Map_{\Sp_{T(n)}^{BS_k}}(\mathbb{S}_{T(n)}\otimes\Sigma^\infty_+S_k,\partial^k(F)) \\
%			\simeq&\Map_{\Fun(\Sp_{T(n)},\Sp_{T(n)})}(\lambda_k(\mathbb{S}_{T(n)}\otimes\Sigma^\infty_+S_k),F) \\
%			\simeq&\Map_{\Fun(\Sp_{T(n)},\Sp_{T(n)})}((-)^{\otimes k},F)
%		\end{align*}
%		where the last equivalence is given by the standard fact that $(\mathbb{S}_{T(n)}\otimes\Sigma^\infty_+S_k\otimes X)_{hS_k}\simeq X$, regardless of the action of $S_k$ on $X$. Then the verification of the $S-k$-action is direct. 
%	\end{proof}

%	\begin{lemma}
%		\[ \partial^{k+1}F\simeq\partial^1(X\mapsto\partial^kF(X\oplus-)) \]
%	\end{lemma}
%	\begin{proof}
%		Note that we have 
%	\end{proof}

%	However, in the $T(n)$-local settings, we have another characterization of the above functors. 

%	\begin{defn}\label{Analytic_fun_defn_2}
%		Consider a finite set $S$, and let $P(S)$ be the poset of subsetes of $S$ ordered by inclusion. We define $P_{\leq i}(S)$ to be the subset of $P(S)$ consisting of subsets having cardinality smaller or equal than $i$ and $P_{>i}(S)$ to be the subset of $P(S)$ consisting of subsets having cardinality greater than $i$. \\
%		We define a $n$-cube $F:P(S)\to\mathcal{C}$ (where $|S|=n$) to be strongly coCartesian if it is the left Kan extension of the subdiagram $F:P_{\leq 1}(S)\to\mathcal{C}$. We define a cube to be Cartesian if it is the right Kan extension of the subdiagram $F:P_{>0}(S)\to\mathcal{C}$. \\
%		We define a functor $F:\mathcal{C}\to\mathcal{D}$ to be $n$-excisive if it takes strongly coCartesian $n$-cubes to Cartesian $n$-cubes. We define the full subcategory $\Exc^n(\mathcal{C},\mathcal{D})$ to be the full subcategory of $\Fun(\mathcal{C},\mathcal{D})$ spanned by the $n$-excisive functors. The inclusion $\Exc^n(\mathcal{C},\mathcal{D})\inj\Fun(\mathcal{C},\mathcal{D})$ admit a left adjoint $P_n$. Moreover, we say $F$ is $n$-homogeneous if $P_{n-1}F\simeq*$ and $P_nF=F$. \\
%		We define a functor $F:\mathcal{C}\to\mathcal{D}$ to be analytic if we have 
%		\[ F\to\plim(P_nF\to\cdots\to P_1F\to P_0F) \]
%		an equivalence. 
%	\end{defn}

%	We need a few inputs to relate the above definitions. 

%	With the above theorems, we can prove the following. 
%	\begin{thm}\label{Ana_fun_tel_loc}
%		The functors introduced in \cref{Analytic_fun_defn_1} are precisely the filtered colimit preserving analytic functors with domain and codomain $\Sp_{T(n)}$, in the sense of \cref{Analytic_fun_defn_2}. 
%	\end{thm}
%	\begin{proof}
%		We observe that the filtered colimit preserving $n$-homogeneous functors are of the form $(C\otimes X^{\otimes n})_{hS_n}$ since any such functor $F:\Sp_{T(n)}^{(n)}\to\Sp_{T(n)}^{(n)}$ is of the form 
%		\[ F(X)=F'(X,\cdots,X)_{hS_n}=(F'(\mathbb{S},\cdots,\mathbb{S})\otimes X^{\otimes n})_{hS_n} \]
%		by \cref{N_homog_fun_char}. This directly proves $\implies$. \\
%		For the converse, an analytic functor $F$ is the colimit of the tower 
%		\[ F=\plim(P_nF\to\cdots\to P_1F\to P_0F) \]
%		Then consider the fiber sequences $D_nF\to P_nF\to P_{n-1}F$. Then to show this fiber sequence splits, it suffices to show that any natural transformation $P_{n-1}F\to\Sigma D_nF$ vanishes. By \cref{N_exc_fun_char}, this is equivalent to that any natural transformation $P_{n-1}F\to(D_nF'\delta)^{tS_n}$ vanishes. But since the codomain is $\Sp_{T(n)}$, in which Tate construction vanishes by \cref{Tate_vanish}, hence this fiber sequence splits. Therefore $P_nF\simeq D_nF\oplus P_{n-1}F$, and hence any analytic functor has the form in \cref{Analytic_fun_defn_1}. 
%	\end{proof}

%	\begin{lemma}\label{Poly_fun_homog_deg_null}
%		For $F,G:\Sp_{T(n)}\to\Sp_{T(n)}$ with $F$ $n$-homogeneous and $G$ $m$-homogeneous ($n\neq m$), any natural transformation $F\to G$ is null-homotopic (that is, factors through the zero map). 
%	\end{lemma}
%	\begin{proof}
%		If $n>m$, then by the definition $P_n$, we have 
%		\[ \Hom_{\Fun(\Sp_{T(n)},\Sp_{T(n)})}(F,G)\simeq\Hom_{\Fun(\Sp_{T(n)},\Sp_{T(n)})}(P_mF,G) \]
%		but by definition $P_mF$ is trivial, hence $F\to G$ must also be trivial. \\
%		If $n<m$, then as in the proof of \cref{Ana_fun_tel_loc} we have any natural transformation $F\to G$ is trivial by Tate vanishing. 
%	\end{proof}

	

%	\begin{coro}\label{Ana_fun_hom_space}
%		We have a fully faithful embedding 
%		\[ \Fun(\Core\Fin,\Sp_{T(n)})\to\Fun(\Sp_{T(n)},\Sp_{T(n)}) \]
%		whose essential image consists of filtered colimit preserving analytic functors. 
%	\end{coro}
%	\begin{proof}
%		By \cref{Ana_fun_tel_loc} and \cref{Poly_fun_homog_deg_null}, it remains to prove that the only natural transformations between functors $F(X)=(C\otimes X^{\otimes m})_{hS_m}$ and $G(X)=(D\otimes X^{\otimes m})_{hS_m}$ comes from a map $C\to D$. Firstly for the case $m=1$, note that 
%	\end{proof}

	%	It suffices to verify that the construction 
%	\[ \lambda_m:\Sp_{T(n)}^{BS_m}\to\Fun(\Sp_{T(n)},\Sp_{T(n)}),C\mapsto(X\mapsto(C\otimes X^{\otimes m})_{hS_m}) \]
%	admits the right adjoint $\partial_m$ sending a functor $F=\bigoplus(C(k)\otimes X^{\otimes k})_{hS_k}$ to the $S_m$ equivariant $T(n)$-local spectrum $C(m)$. Firstly clearly $\lambda_m$ preserves colimits, it admit a right adjoint. 